{\centering
	\chapter{Construcción}\label{parte2}}
\thispagestyle{empty}
\noindent
\begin{minipage}[t]{0.35\textwidth}
	\begin{flushright}
		\large
		\vspace{1pt}
		三十輻共一轂，\\
		當其無，有車之用。\\
		\vspace{0.5em}
		埏埴以為器，\\
		當其無，有器之用。\\
		\vspace{0.5em}
		鑿戶牖以為室，\\
		當其無，有室之用。\\
		\vspace{0.5em}
		故有之以為利，\\
		無之以為用。
	\end{flushright}
\end{minipage}%
\hfill
\vrule width 0.5pt % La línea vertical
\hfill
\begin{minipage}[t]{0.55\textwidth}
	\vspace{1pt}
	\raggedright
	\emph
	Treinta radios convergen en el centro de una rueda,\\
	pero es su vacío lo que hace útil al carro.\\
	
	\vspace{0.8em}
	Se moldea la arcilla para hacer un recipiente,\\
	pero es en su vacío donde está la utilidad del recipiente.\\
	
	\vspace{0.8em}
	Se abren puertas y ventanas en los muros de una casa,\\
	y es por sus vacíos que la casa es habitable.\\
	
	\vspace{0.8em}
	En el ser centramos nuestro interés,\\
	pero del no-ser depende la utilidad.
\end{minipage}

\vspace{1em}
\hfill {\normalfont  道德經 —\textsc{Dàodé Jīng, Capítulo 11}}

\section{El Go en dos dimensiones y dos jugadores.}

Interpretemos el tablero de forma conveniente, consideremos únicamente sus puntos de intersección. Sean $  n, m \in \mathbb{N} $. Definimos $ N = \{1, \dots, n\}  $, $  M = \{1, \dots, m\}$, con estos conjuntos, formalizamos el espacio de juego,\footnote{Utilizamos la letra $G$ porque emplearemos la terminología japonesa del Go para la notación, siempre que no resulte confuso. En dicho idioma, el tablero se denomina \textit{Goban} (碁盤).}
$  G \coloneqq N \times M = \{(x,y) : x \in N,\, y \in M\}.  $
Sean $t_1, t_2 \in G$ dos \textit{puntos}\footnote{Se usa la letra $t$ por el término \textit{ten} (点), punto en japonés.}. Si $t_1=(a,b)$, diremos que $t_2$ es adyacente a $t_1$ y lo denotaremos como $t_1 \sim t_2$ si cumple exactamente una de las siguientes condiciones,
$$t_2 = (a+1,b), \quad t_2 = (a,b+1), \quad t_2 = (a-1,b), \quad t_2 = (a,b-1).$$
Si no se da ninguno de estos casos, $t_{1}$ no es adyacente a $t_{2}$, lo cual denotaremos como $t_{1} \nsim t_{2}$. La Figura~\ref{fig:Fig 13} ilustra las adyacencias de un punto.

\begin{figure}[H]
	\centering
	\begin{tikzpicture}[scale=0.485] \tikzstyle{every node}=[font=\small]
		% Draw the background
		\fill[goban] (-0.5,-0.5) rectangle (6.5,6.5);
		
		% Draw the grid lines
		\draw[step=1cm, black] (0,0) grid (6,6);
		
		% Draw the thick border
		\draw[very thick] (0,0) rectangle (6,6);
		
		% Draw the central black dot
		\filldraw[black] (3,3) circle (4pt);
		
		\fill[Blue] (3.75,3.75) rectangle (4.25,4.25);
		\fill[OliveGreen] (3,4) circle (5pt);
		\fill[OliveGreen] (5,4) circle (5pt);
		\fill[OliveGreen] (4,3) circle (5pt);
		\fill[OliveGreen] (4,5) circle (5pt);
	\end{tikzpicture}
	\caption{Punto $(5,5)$ en \textcolor{Blue}{azul}, adyacencias en \textcolor{OliveGreen}{verde}.}
	\label{fig:Fig 13}
\end{figure}

Teniendo el \textit{dónde}, veamos \textit{quiénes} participan en él. Cuando vimos las reglas, notamos que cada punto del tablero puede estar en uno de tres estados posibles: no tener piedra, tener una piedra blanca o tener una piedra negra. Definimos $E$ como el conjunto de estados que un punto puede adoptar:
$$E=\{0,1,2\}.$$
Estos valores representan, respectivamente: libre, negro y blanco. Con esto, podemos dar cuenta de todos los estados posibles en el tablero:
\begin{align*}
	E \times G &= \{(e,t) : e \in E, t \in G \}.
\end{align*}
Por ejemplo, el punto $(4,4) \in G$ puede encontrarse en uno de tres estados: $(0,(4,4))$, sin piedra en $(4,4)$; $(1,(4,4))$, piedra negra en $(4,4)$; o $(2,(4,4))$, piedra blanca en $(4,4)$.

A los elementos de $E \times G$ los llamaremos \textit{piedras} y los denotaremos con la letra $i$, de modo que $i=(e,t)$\footnote{Por el término \textit{ishi} (石), que significa «piedra» en japonés.}, con $e \in E$ y $t \in G$. Pero, en el Go, un mismo punto no puede tener todos los estados al mismo tiempo, a lo sumo tendrá uno.
\begin{definition}[Tablero]
	Un subconjunto $T$ de $E\times G$ se dirá un \textit{tablero} si, dado $(e,t)\in T$, $$(e',t)\notin T, \forall e' \neq e.$$
\end{definition}
Es decir, en $T$ hay un único estado para cada posición. Teniendo \textit{dónde} y \textit{quiénes}, podemos construir el \textit{qué}. Para ello necesitamos de dos funciones esenciales.
\begin{definition}[Función estado de una piedra]\label{funcioncolor}
	\[\text{Definimos } C: E \times G \to E, \quad (e,t) \mapsto e.\]
\end{definition}
\begin{definition}[Función posición de una piedra]\label{funcionposiciion}
	\[\text{Definimos } P: E \times G \to G, \quad (e,t) \mapsto t.\]
\end{definition}
Con ellas, el concepto de grupo emerge naturalmente.
\begin{definition}[Conexión]
	Sea $T$ un tablero e $i_1, i_2 \in T$ dos piedras. Se dice que $i_1$ e $i_2$ están conectadas y lo notaremos por $i_1 \leftrightarrow i_2$ si,
	\[C(i_1) = C(i_2) \text{ y } P(i_1) \sim P(i_2).\]
\end{definition}
\begin{definition}[Grupo de piedras]
	Dado un tablero $T$, dos piedras $i_a$ e $i_b$ pertenecen al mismo grupo (lo cual denotaremos como $i_a \equiv i_b$) si existe una secuencia finita de piedras $\{i_1, i_2, \dots, i_k\}$ tal que $i_1 = i_a$, $i_k = i_b$, y cada piedra está conectada con la siguiente ($i_m \leftrightarrow i_{m+1}$). La relación $\equiv$ es de equivalencia.
	
	Notaremos a los grupos como $I^e$, donde $e$ indica el estado o color del grupo, o diremos «\textit{Sea $I$ un grupo de color $e$}» o bien simplemente $I$ si el color no resulta relevante.\footnote{\textit{ishi} (石, piedra/s). Igual que las piedras, pero en mayúscula al tratarse de un conjunto de ellas.}
\end{definition}

\begin{definition}[Adyacencias de un grupo]
	Sea $I$ un grupo de color $e$, $i=(e',t)$ una piedra del tablero. $I$ es adyacente a $i$ y lo denotaremos como $I \sim i$ si al menos una piedra del grupo es adyacente a $i$:
	\[I \sim i \iff \exists \ i_k \in I : P(i_k) \sim P(i).\]
\end{definition}

Estas definiciones abarcan también a los grupos $I^{0}$. Dado que estos grupos son de vital importancia, recibirán un nombre particular: los llamaremos \textit{regiones} y los denotaremos simplemente con la letra $R$.

Para referirnos al color de un grupo en su totalidad, definimos $C(I^{e}) = e$, basándonos en que todos los elementos que lo componen comparten el mismo estado.

\begin{wrapfigure}[13]{r}{0.35\textwidth}
	%\begin{figure}[H]
	\centering
	\begin{tikzpicture}[scale=0.485] \tikzstyle{every node}=[font=\small]
		% Draw the background
		\fill[goban] (-0.5,-0.5) rectangle (6.5,6.5);
		
		% Draw the grid lines
		\draw[step=1cm, black] (0,0) grid (6,6);
		
		% Draw the thick border
		\draw[very thick] (0,0) rectangle (6,6);
		
		% Draw the central black dot
		
		\filldraw[black] (2,4) circle (12pt);
		\filldraw[black] (3,4) circle (12pt);
		\filldraw[black] (4,3) circle (12pt);
		\filldraw[black] (4,4) circle (12pt);
		\filldraw[draw=black,fill=white] (3,3) circle (12pt);
		\filldraw[draw=black,fill=white] (3,2) circle (12pt);
		\fill[OliveGreen] (3,3) circle (5pt);
		\fill[OliveGreen] (1,4) circle (5pt);
		\fill[OliveGreen] (5,4) circle (5pt);
		\fill[OliveGreen] (3,5) circle (5pt);
		\fill[OliveGreen] (2,5) circle (5pt);
		\fill[OliveGreen] (4,5) circle (5pt);
		\fill[OliveGreen] (4,2) circle (5pt);
		\fill[OliveGreen] (2,3) circle (5pt);
		\fill[OliveGreen] (5,3) circle (5pt);
		
	\end{tikzpicture}
	\caption{Grupo de piedras negras con sus \textcolor{OliveGreen}{adyacencias} marcadas.}
	\label{ejemplo_adyacencias}
\end{wrapfigure}Nos resultará muy útil conocer las adyacencias de una piedra o grupo. Para ello, definiremos el conjunto $K$\footnote{Por el término \textit{chikaku} (近く,cerca).} asociado a una piedra o grupo. véase Figura~\ref{ejemplo_adyacencias}



\begin{definition}[Conjunto de adyacencias]
	Sea $i$ una piedra e $I$ un grupo de piedras de color $e$. Definimos el conjunto de adyacencias del grupo como las piedras que cumplen la definición de adyacencia, esto es:
	\[K(I) = \{i_r : I \sim i_r\}.\]
	Análogamente, lo definimos para una piedra $i$:
	\[K(i) = \{i_r : P(i) \sim P(i_r)\}.\]
\end{definition}

Observemos que toda piedra adyacente del mismo color ya pertenece al grupo. Así, si el grupo es de color $e$, los estados de las adyacencias pertenecerán, a lo sumo, a $E\setminus\{e\}$.

Para dar cuenta de los estados de las adyacencias, haremos uso (y abuso) de las funciones que hemos definido.

\begin{definition}[Estado de las adyacencias de un grupo o piedra]\label{def_estados}
	Sea $I$ un grupo de color $e$. Al aplicar la función $C$ elemento a elemento sobre el conjunto de sus adyacencias $K(I)$, obtenemos el conjunto que buscamos. Lo definimos como:
	\[C(K(I)) = \{C(i) \mid   i\in K(I)\} \subset \{0, e'\}.\]
	Análogamente, para una piedra individual $i = (e,t)$, lo definimos como:
	\[C(K(i)) = \{C(i_k) \mid i_k \in K(i)\} \subset \{0, e'\}.\]
\end{definition}
Esta definición es fundamental para describir gran parte de la dinámica del Go. De ahora en adelante, siempre que hablemos de grupos $I$, ha de interpretarse que son de color distinto de $0$. Siempre que tratemos con $I^{0}$, utilizaremos la notación $R$, salvo que se indique lo contrario.

\begin{wrapfigure}[11]{r}{0.35\textwidth}
	\vspace{-36pt}
	\centering
	\begin{tikzpicture}[scale=0.485] \tikzstyle{every node}=[font=\small]
		% Draw the background
		\fill[goban] (-0.5,-0.5) rectangle (6.5,6.5);
		
		% Draw the grid lines
		\draw[step=1cm, black] (0,0) grid (6,6);
		
		% Draw the thick border
		\draw[very thick] (0,0) rectangle (6,6);
		
		% Draw the central black dot
		\filldraw[draw=black,fill=white] (2,3) circle (12pt);
		\filldraw[draw=black,fill=white] (1,3) circle (12pt);
		\filldraw[draw=black,fill=white] (3,3) circle (12pt);
		\filldraw[draw=black,fill=white] (1,5) circle (12pt);
		\filldraw[draw=black,fill=white] (1,4) circle (12pt);
		\filldraw[draw=black,fill=white] (4,2) circle (12pt);
		\filldraw[draw=black,fill=white] (4,1) circle (12pt);
		\filldraw[draw=black,fill=white] (4,0) circle (12pt);
		\filldraw[draw=black,fill=white] (0,5) circle (12pt);
		\filldraw[black] (5,3) circle (12pt);
		\filldraw[black] (5,4) circle (12pt);
		\filldraw[black] (4,4) circle (12pt);
		\filldraw[black] (3,4) circle (12pt);
		\filldraw[black] (2,5) circle (12pt);
		\filldraw[black] (2,6) circle (12pt);
		\filldraw[black] (5,2) circle (12pt);
		\filldraw[black] (5,1) circle (12pt);
		\filldraw[black] (5,0) circle (12pt);
		
		\filldraw[Blue] (0,0) circle (6pt);
		\filldraw[Blue] (0,1) circle (6pt);
		\filldraw[Blue] (0,2) circle (6pt);
		\filldraw[Blue] (0,3) circle (6pt);
		\filldraw[Blue] (0,4) circle (6pt);
		\filldraw[Blue] (1,0) circle (6pt);
		\filldraw[Blue] (1,1) circle (6pt);
		\filldraw[Blue] (1,2) circle (6pt);
		\filldraw[Blue] (2,0) circle (6pt);
		\filldraw[Blue] (2,1) circle (6pt);
		\filldraw[Blue] (2,2) circle (6pt);
		\filldraw[Blue] (3,0) circle (6pt);
		\filldraw[Blue] (3,1) circle (6pt);
		\filldraw[Blue] (3,2) circle (6pt);
		
		\filldraw[OliveGreen] (6,6) circle (6pt);
		\filldraw[OliveGreen] (6,5) circle (6pt);
		\filldraw[OliveGreen] (6,4) circle (6pt);
		\filldraw[OliveGreen] (6,3) circle (6pt);
		\filldraw[OliveGreen] (6,2) circle (6pt);
		\filldraw[OliveGreen] (6,1) circle (6pt);
		\filldraw[OliveGreen] (6,0) circle (6pt);
		\filldraw[OliveGreen] (5,6) circle (6pt);
		\filldraw[OliveGreen] (4,6) circle (6pt);
		\filldraw[OliveGreen] (3,6) circle (6pt);
		\filldraw[OliveGreen] (5,6) circle (6pt);
		\filldraw[OliveGreen] (4,6) circle (6pt);
		\filldraw[OliveGreen] (3,6) circle (6pt);
		\filldraw[OliveGreen] (5,5) circle (6pt);
		\filldraw[OliveGreen] (4,5) circle (6pt);
		\filldraw[OliveGreen] (3,5) circle (6pt);
		\filldraw[OliveGreen] (5,6) circle (6pt);
		\filldraw[OliveGreen] (4,6) circle (6pt);
		\filldraw[Orange] (0,6) circle (6pt);
		\filldraw[Orange] (1,6) circle (6pt);
		\filldraw[Orange] (2,4) circle (6pt);
		\filldraw[Orange] (4,3) circle (6pt);
		\filldraw[Blue] (0,0) circle (6pt);
	\end{tikzpicture}
	\caption{Territorio \textcolor{Blue}{blanco}, \textcolor{OliveGreen}{negro} y \textcolor{Orange}{libre}.}
	\label{fig:territorios}
\end{wrapfigure}

\begin{definition}[Territorio]
	Una región $R$ es territorio del jugador de color $e$ si existe un conjunto $S$ de grupos de color $e$ tal que,
	\begin{align*}
		K(R) \subseteq \bigcup_{I\in S} I.
	\end{align*}
\end{definition}
Es decir, todas las piedras adyacentes a la región pertenecen a grupos del mismo color. Denotaremos los \textit{territorios}\footnote{Por el término \textit{ji} (地), que significa «tierra» en japonés.} de cada color como,
\vspace{-1em}
$$J^{e} = \{R : C(K(R)) = \{e\}\}, e\neq 0.$$

A las regiones que no son territorio de ningún jugador las notaremos y definiremos como,
\vspace{-1em}
$$J^{0}\coloneqq \{R : |C(K(R))| \geq 2\}.$$
La Figura~\ref{fig:territorios} ilustra los diferentes tipos de \textit{territorios}.


Con esto, hemos construido una partición del conjunto de \textit{regiones} ($I^{0}$), ya que toda región pertenece al territorio blanco o pertenece al territorio negro o bien, no forma parte de ningún territorio:
\begin{align*}
	I^{0} = J^{0} \cup J^{1} \cup J^{2}.
\end{align*}
\label{posicion_espacial}
Adicionalmente, la cardinalidad del conjunto de adyacencias de una piedra individual, $|K(i)|$, es de gran utilidad para determinar su ubicación espacial en el tablero. Dado que $|K(i)|$ se encuentra acotado por el Lema~\ref{lem:2n}, un simple análisis de la cantidad de adyacencias nos revela exactamente si la piedra está situada en el interior, en un borde o en una esquina. Esta propiedad geométrica será esencial y haremos uso de ella en la Sección~\ref{teoremas_go1} y \ref{sec_JR}.

Para formalizar las capturas, introducimos primero el concepto de \textit{libertad}.

\begin{definition}[Libertades de un grupo]
	Sea $I$ un grupo. Definimos el conjunto de libertades de dicho grupo como:
	\[L(I) = \{i \in K(I) : C(i) = 0\}.\]
\end{definition}
Notemos que, si para un grupo ocurre que $L(I) = \varnothing$, entonces dicho grupo se considera \textit{muerto}. Esta situación solo puede darse al momento de realizar una jugada y concretar una captura. Pero, ¿qué es una captura? ¿qué es una jugada?


Un grupo es capturado cuando pierde todas sus libertades. Las piedras capturadas se retiran del tablero\footnote{Nosotros hacemos uso de las reglas chinas para contar los puntos, si se juega con otro método, japonesas por ejemplo, habría que otorgar un punto por piedra capturada.}.

\begin{definition}[Criterio de captura]
	Sea $I$ un grupo, entonces,
	\[I \text{ es capturado} \iff L(I) = \varnothing.\]
	Para formalizar el retiro de estas piedras del tablero, si $I = \{i_1, i_2, \dots, i_n\}$, basta reasignar sus estados haciendo $C(i_j) = 0$ para todo $j=1, \dots, n$.
\end{definition}
Preferimos el término \textit{es capturado}, ya que la muerte es consecuencia de una captura.

\section{Hacia el movimiento}

Hasta el momento hemos descrito la estructura estática del tablero. Para modelar una partida debemos incorporar la evolución temporal del juego. Para ello, definimos primero el espacio de todas las jugadas posibles, al que denotaremos como $\mathcal{J}$,
$$ \mathcal{J} \coloneqq \left((E \setminus \{0\}) \times G\right) \cup \{\varnothing\}.$$

\subsection{Función de jugada}
\label{criterio_jugada}
Al jugar al Go y ser nuestro turno, tomamos una piedra y la colocamos en el tablero o bien pasamos. Definimos la función de jugada como, $f: \mathbb{N} \to \mathcal{J}$. Veamos un ejemplo.
\begin{wrapfigure}[4]{r}{0.4\textwidth}
	%\begin{figure}[H]
	\centering
	\begin{tikzpicture}[scale=0.485] \tikzstyle{every node}=[font=\small]
		
		% Draw the background
		\fill[goban] (-0.5,-0.5) rectangle (6.5,6.5);
		
		% Draw the grid lines
		\draw[step=1cm, black] (0,0) grid (6,6);
		
		% Draw the thick border
		\draw[very thick] (0,0) rectangle (6,6);
		
		% Draw the central black dot
		\filldraw[black] (2,2) circle  (4pt) ;
		\path (2,2) node [white]{1} ;
		
		\filldraw[black] (2,2) circle (12pt);
		\path (2,2) node [white]{1} ;
		
		\filldraw[black] (2,3) circle (12pt);
		\path (2,3) node [white]{3} ;
		
		\filldraw[black] (3,3) circle (12pt);
		\path (3,3) node [white]{5} ;
		
		\filldraw[black] (4,5) circle (12pt);
		\path (4,5) node [white]{7} ;
		
		\filldraw[black] (3,5) circle (12pt);
		\path (3,5) node [white]{9} ;
		
		\filldraw[draw=black,fill=white] (4,4) circle (12pt);
		\path (4,4) node [black]{2} ;
		\filldraw[draw=black,fill=white] (4,2) circle (12pt);
		\path (4,2) node [black]{4} ;
		\filldraw[draw=black,fill=white] (3,2) circle (12pt);
		\path (3,2) node [black]{6} ;
		\filldraw[draw=black,fill=white] (4,3) circle (12pt);
		\path (4,3) node [black]{8} ;
		\filldraw[draw=black,fill=white] (2,1) circle (12pt);
		\path (2,1) node [black]{10} ;
		\filldraw[draw=black,fill=white] (1,2) circle; 
	\end{tikzpicture}
\end{wrapfigure} \begin{align*}
	f(1) &= (1,(2,2)), & f(2) &= (1,(4,4)), & f(3) &= (1,(2,3)),\\
	f(4) &= (2,(4,2)), & f(5) &= (2,(3,3)), & f(6) &= (2,(3,2)),\\
	f(7) &= (1,(4,5)), & f(8) &= (1,(3,4)), & f(9) &= (1,(3,5)),\\
	f(10) &= (2,(2,1)), & f(11) &= \varnothing, & f(12) &= \varnothing.
\end{align*}
%\vspace{2.5cm} 
\break

\subsection{Reglas de la jugada legal}
\begin{wrapfigure}[12]{r}{0.5\textwidth}
	\vspace{-15pt}
	%\begin{figure}[H]
	\centering
	\begin{tikzpicture}[scale=0.485] \tikzstyle{every node}=[font=\small]
		% Draw the background
		\fill[goban] (-0.5,-0.5) rectangle (6.5,6.5);
		% Draw the grid lines
		\draw[step=1cm, black] (0,0) grid (6,6);
		% Draw the thick border
		\draw[very thick] (0,0) rectangle (6,6);
		% Draw the central black dot
		\filldraw[black] (3,3) circle (4pt);
		
		\filldraw[draw=black,fill=black] (4,4) circle (12pt);
		\filldraw[draw=black,fill=black] (4,2) circle (12pt);
		\filldraw[draw=black,fill=black] (3,2) circle (12pt);
		\filldraw[draw=black,fill=black] (3,3) circle (12pt);
		\filldraw[draw=black,fill=black] (5,2) circle (12pt);
		\filldraw[draw=black,fill=black] (6,3) circle (12pt);
		\filldraw[draw=black,fill=white] (5,3) circle (12pt);
		\filldraw[draw=black,fill=black] (5,4) circle (12pt);
		\fill[draw=black,fill=BlueViolet, very thin] (3.75,2.75) rectangle (4.25,3.25);
		\path (3,7.25) node [black,very thick]{\large{Juegan Blancas}} ;
		
		
		\begin{scope}[xshift=8cm]
			% Draw the background
			\fill[goban] (-0.5,-0.5) rectangle (6.5,6.5);
			% Draw the grid lines
			\draw[step=1cm, black] (0,0) grid (6,6);
			% Draw the thick border
			\draw[very thick] (0,0) rectangle (6,6);
			% Draw the central black dot
			\filldraw[black] (3,3) circle (4pt);
			
			\filldraw[draw=black,fill=white] (4,4) circle (12pt);
			\filldraw[draw=black,fill=white] (4,2) circle (12pt);
			\filldraw[draw=black,fill=white] (3,2) circle (12pt);
			\filldraw[draw=black,fill=white] (3,3) circle (12pt);
			\filldraw[draw=black,fill=white] (5,2) circle (12pt);
			\filldraw[draw=black,fill=white] (6,3) circle (12pt);
			\filldraw[black] (5,3) circle (12pt);
			\fill[draw=black,fill=BlueViolet, very thin] (3.75,2.75) rectangle (4.25,3.25);
			\filldraw[draw=black,fill=white] (5,4) circle (12pt);
			\path (3,7.25) node [black,very thick]{\large{Juegan Negras}} ;
		\end{scope}
		
	\end{tikzpicture}
	\caption{En \textcolor{BlueViolet}{violeta} las jugadas ilegales, pues llevan al suicidio.}
	\label{sui_ej}
\end{wrapfigure}

La función registra las jugadas, pero no su legalidad. A continuación, analizaremos estas reglas. La \textit{regla del suicidio} prohíbe colocar una piedra si esto provoca la captura inmediata del propio grupo, Figura~\ref{sui_ej}. Existe una única excepción: una jugada sin libertades es legal si captura simultáneamente un grupo de otro color, véase la Figura~\ref{fig: fig1.19}. Notemos que, debido a esta mecánica, cuando es el turno del negro solo pueden ocurrir capturas de grupos blancos, y viceversa. \break 

\begin{wrapfigure}[11]{r}{0.5\textwidth}
	
	%\begin{figure}[H]
	\centering
	\begin{tikzpicture}[scale=0.485] \tikzstyle{every node}=[font=\small]
		% Draw the background
		\fill[goban] (-0.5,-0.5) rectangle (6.5,6.5);
		
		% Draw the grid lines
		\draw[step=1cm, black] (0,0) grid (6,6);
		
		% Draw the thick border
		\draw[very thick] (0,0) rectangle (6,6);
		
		% Draw the central black dot
		\filldraw[draw=black,fill=white] (3,4) circle (12pt);
		\filldraw[draw=black,fill=white] (2,3) circle (12pt);
		\filldraw[draw=black,fill=white] (4,3) circle (12pt);
		\filldraw[draw=black,fill=white] (3,2) circle (12pt);
		\filldraw[black] (3,1) circle (12pt);
		\filldraw[black] (2,2) circle (12pt);
		\filldraw[black] (4,2) circle (12pt);
		\path (3,2) node [black]{1};
		\begin{scope}[xshift=8cm]
			% Draw the background
			\fill[goban] (-0.5,-0.5) rectangle (6.5,6.5);
			
			% Draw the grid lines
			\draw[step=1cm, black] (0,0) grid (6,6);
			
			% Draw the thick border
			\draw[very thick] (0,0) rectangle (6,6);
			
			% Draw the central black dot
			\filldraw[draw=black,fill=white] (3,4) circle (12pt);
			\filldraw[draw=black,fill=white] (2,3) circle (12pt);
			\filldraw[draw=black,fill=white] (4,3) circle (12pt);
			\filldraw[black] (3,3) circle (12pt);
			\filldraw[black] (3,1) circle (12pt);
			\filldraw[black] (2,2) circle (12pt);
			\filldraw[black] (4,2) circle (12pt);
			\path (3,3) node [white]{2};
		\end{scope}
	\end{tikzpicture}
	\caption{Posición de una piedra bajo suicidio ya que captura.}
	\label{fig: fig1.19}
\end{wrapfigure} Para evaluar una jugada, primero verificamos si produce una captura. En ese caso, la jugada es legal. Si no captura, comprobamos sus libertades: si no tiene ninguna, es un suicidio y, por tanto, ilegal; en caso contrario, es legal. Como mencionamos en la Sección~\ref{primerreglako}, debemos tener presente la \textit{regla del ko}. Esta prohíbe la recaptura inmediata de una única piedra si dicha acción revierte el estado del tablero exactamente a la posición del turno previo véase la Figura~\ref{regla_ko_ej}. \break

\vspace{15pt}
\begin{wrapfigure}[6]{r}{0.5\textwidth}
	\vspace{-10pt}
	%\begin{figure}[H]
	\centering
	\begin{tikzpicture}[scale=0.485] \tikzstyle{every node}=[font=\small]
		% Draw the background
		\fill[goban] (-0.5,-0.5) rectangle (6.5,6.5);
		
		% Draw the grid lines
		\draw[step=1cm, black] (0,0) grid (6,6);
		
		% Draw the thick border
		\draw[very thick] (0,0) rectangle (6,6);
		
		% Draw the central black dot
		\filldraw[draw=black,fill=white] (3,4) circle (12pt);
		\filldraw[draw=black,fill=white] (2,3) circle (12pt);
		\filldraw[draw=black,fill=white] (4,3) circle (12pt);
		\filldraw[draw=black,fill=white] (3,2) circle (12pt);
		\filldraw[black] (3,1) circle (12pt);
		\filldraw[black] (2,2) circle (12pt);
		\filldraw[black] (4,2) circle (12pt);
		\path (3,2) node [black]{1};
		\begin{scope}[xshift=8cm]
			% Draw the background
			\fill[goban] (-0.5,-0.5) rectangle (6.5,6.5);
			
			% Draw the grid lines
			\draw[step=1cm, black] (0,0) grid (6,6);
			
			% Draw the thick border
			\draw[very thick] (0,0) rectangle (6,6);
			
			% Draw the central black dot
			\filldraw[draw=black,fill=white] (3,4) circle (12pt);
			\filldraw[draw=black,fill=white] (2,3) circle (12pt);
			\filldraw[draw=black,fill=white] (4,3) circle (12pt);
			\filldraw[black] (3,3) circle (12pt);
			\filldraw[OliveGreen] (3,2) circle (6pt);  
			\filldraw[black] (3,1) circle (12pt);
			\filldraw[black] (2,2) circle (12pt);
			\filldraw[black] (4,2) circle (12pt);
			\path (3,3) node [white]{2};
		\end{scope}
	\end{tikzpicture}
	\caption{Blanco no puede jugar en el \textcolor{OliveGreen}{punto verde}, pues repite posición.}
	\label{regla_ko_ej}
\end{wrapfigure}La \textit{regla del ko} es insuficiente, pues existen ciclos donde, sin recapturar la piedra inmediata, se repite el tablero. Véase la Figura~\ref{ciclo_ko}, donde se llega a la misma posición que antes, pero a través de capturas que no van en contra de la \textit{regla del ko}.
\vfill
\begin{figure}[H]
	\centering
	\begin{tikzpicture}[scale=0.485] \tikzstyle{every node}=[font=\small]
		% Draw the background
		\fill[goban] (-0.5,-0.5) rectangle (6.5,6.5);
		
		% Draw the grid lines
		\draw[step=1cm, black] (0,0) grid (6,6);
		
		% Draw the thick border
		\draw[very thick] (0,0) rectangle (6,6);
		
		% Draw the central black dot
		\filldraw[draw=black,fill=white] (0,5) circle (12pt);
		\filldraw[black] (1,6) circle (12pt);
		\filldraw[black] (1,5) circle (12pt);
		\filldraw[draw=black,fill=white] (0,3) circle (12pt);
		\filldraw[draw=black,fill=white] (1,3) circle (12pt);
		\filldraw[draw=black,fill=white] (1,4) circle (12pt);
		\filldraw[draw=black,fill=white] (2,4) circle (12pt);
		\filldraw[draw=black,fill=white] (2,5) circle (12pt);
		\filldraw[draw=black,fill=white] (3,5) circle (12pt);
		\filldraw[draw=black,fill=white] (3,6) circle (12pt);
		\filldraw[black] (4,6) circle (12pt);
		\filldraw[black] (4,5) circle (12pt);
		\filldraw[black] (4,4) circle (12pt);
		\filldraw[black] (3,4) circle (12pt);
		\filldraw[black] (2,3) circle (12pt);
		\filldraw[black] (2,2) circle (12pt);
		\filldraw[black] (1,2) circle (12pt);
		\filldraw[black] (0,2) circle (12pt);
		\filldraw[black] (4,2) circle (12pt);
		\begin{scope}[xshift=7.5cm]
			% Draw the background
			\fill[goban] (-0.5,-0.5) rectangle (6.5,6.5);
			
			% Draw the grid lines
			\draw[step=1cm, black] (0,0) grid (6,6);
			
			% Draw the thick border
			\draw[very thick] (0,0) rectangle (6,6);
			
			% Draw the central black dot
			\filldraw[draw=black,fill=white] (0,5) circle (12pt);
			\filldraw[black] (1,6) circle (12pt);
			\filldraw[black] (1,5) circle (12pt);
			\filldraw[draw=black,fill=white] (0,3) circle (12pt);
			\filldraw[draw=black,fill=white] (1,3) circle (12pt);
			\filldraw[draw=black,fill=white] (1,4) circle (12pt);
			\filldraw[draw=black,fill=white] (2,4) circle (12pt);
			\filldraw[draw=black,fill=white] (2,5) circle (12pt);
			\filldraw[draw=black,fill=white] (3,5) circle (12pt);
			\filldraw[draw=black,fill=white] (3,6) circle (12pt);
			\filldraw[black] (4,6) circle (12pt);
			\filldraw[black] (4,5) circle (12pt);
			\filldraw[black] (4,4) circle (12pt);
			\filldraw[black] (3,4) circle (12pt);
			\filldraw[black] (2,3) circle (12pt);
			\filldraw[black] (2,2) circle (12pt);
			\filldraw[black] (1,2) circle (12pt);
			\filldraw[black] (0,2) circle (12pt);
			\filldraw[black] (4,2) circle (12pt);
			\filldraw[draw=black,fill=white] (0,6) circle (12pt);
			\path (0,6) node [black]{1};
		\end{scope}
		\begin{scope}[xshift=15cm]
			% Draw the background
			\fill[goban] (-0.5,-0.5) rectangle (6.5,6.5);
			
			% Draw the grid lines
			\draw[step=1cm, black] (0,0) grid (6,6);
			
			% Draw the thick border
			\draw[very thick] (0,0) rectangle (6,6);
			
			% Draw the central black dot
			\filldraw[black] (1,6) circle (12pt);
			\filldraw[black] (1,5) circle (12pt);
			\filldraw[draw=black,fill=white] (0,3) circle (12pt);
			\filldraw[draw=black,fill=white] (1,3) circle (12pt);
			\filldraw[draw=black,fill=white] (1,4) circle (12pt);
			\filldraw[draw=black,fill=white] (2,4) circle (12pt);
			\filldraw[draw=black,fill=white] (2,5) circle (12pt);
			\filldraw[draw=black,fill=white] (3,5) circle (12pt);
			\filldraw[draw=black,fill=white] (3,6) circle (12pt);
			\filldraw[black] (4,6) circle (12pt);
			\filldraw[black] (4,5) circle (12pt);
			\filldraw[black] (4,4) circle (12pt);
			\filldraw[black] (3,4) circle (12pt);
			\filldraw[black] (2,3) circle (12pt);
			\filldraw[black] (2,2) circle (12pt);
			\filldraw[black] (1,2) circle (12pt);
			\filldraw[black] (0,2) circle (12pt);
			\filldraw[black] (4,2) circle (12pt);
			\filldraw[black] (0,4) circle (12pt);
			\path (0,4) node [white]{2};
		\end{scope}
		
		\begin{scope}[xshift=22.5cm]
			% Draw the background
			\fill[goban] (-0.5,-0.5) rectangle (6.5,6.5);
			
			% Draw the grid lines
			\draw[step=1cm, black] (0,0) grid (6,6);
			
			% Draw the thick border
			\draw[very thick] (0,0) rectangle (6,6);
			
			% Draw the central black dot
			\filldraw[black] (1,6) circle (12pt);
			\filldraw[black] (1,5) circle (12pt);
			\filldraw[draw=black,fill=white] (0,3) circle (12pt);
			\filldraw[draw=black,fill=white] (1,3) circle (12pt);
			\filldraw[draw=black,fill=white] (1,4) circle (12pt);
			\filldraw[draw=black,fill=white] (2,4) circle (12pt);
			\filldraw[draw=black,fill=white] (2,5) circle (12pt);
			\filldraw[draw=black,fill=white] (3,5) circle (12pt);
			\filldraw[draw=black,fill=white] (3,6) circle (12pt);
			\filldraw[black] (4,6) circle (12pt);
			\filldraw[black] (4,5) circle (12pt);
			\filldraw[black] (4,4) circle (12pt);
			\filldraw[black] (3,4) circle (12pt);
			\filldraw[black] (2,3) circle (12pt);
			\filldraw[black] (2,2) circle (12pt);
			\filldraw[black] (1,2) circle (12pt);
			\filldraw[black] (0,2) circle (12pt);
			\filldraw[black] (4,2) circle (12pt);
			\filldraw[draw=black,fill=white] (0,5) circle (12pt);
			\path (0,5) node [black]{3};
		\end{scope}
	\end{tikzpicture}
	\caption{Tablero con un ciclo de tres jugadas.}
	\label{ciclo_ko}
\end{figure}
\newpage

Como la repetición de posiciones está prohibida, necesitamos una regla más general. Para ilustrar un caso extremo, véase la Figura~\ref{ciclo_inf}.\footnote{Existen ejemplos con ciclos mucho más grandes; para profundizar en ellos, véase \cite{cicloinfinito}.}
\begin{figure}[H]
	\centering
	\begin{tikzpicture}[scale=0.75]
		% Draw the background
		\fill[goban] (-0.5,-0.5) rectangle (2.5,2.5);
		% Draw the grid lines
		\draw[step=1cm, black] (0,0) grid (2,2);
		% Draw the thick border
		\draw[very thick] (0,0) rectangle (2,2);
		
		% Draw the central black dot
		\filldraw[draw=black,fill=white] (0,2) circle (12pt);
		\filldraw[draw=black,fill=white] (1,1) circle (12pt);
		\filldraw[draw=black,fill=white] (0,0) circle (12pt);
		\filldraw[black] (1,2) circle (12pt);
		\filldraw[black] (2,1) circle (12pt);
		\filldraw[black] (1,0) circle (12pt);
		
		\begin{scope}[xshift=3.5cm]
			% Draw the background
			\fill[goban] (-0.5,-0.5) rectangle (2.5,2.5);
			% Draw the grid lines
			\draw[step=1cm, black] (0,0) grid (2,2);
			% Draw the thick border
			\draw[very thick] (0,0) rectangle (2,2);
			\filldraw[draw=black,fill=white] (0,2) circle (12pt);
			\filldraw[draw=black,fill=white] (1,1) circle (12pt);
			\filldraw[draw=black,fill=white] (0,0) circle (12pt);
			\filldraw[black] (2,1) circle (12pt);
			\filldraw[black] (1,0) circle (12pt);
			\filldraw[draw=black,fill=white] (2,2) circle (12pt);
			\path (2,2) node [black]{1};
		\end{scope}
		\begin{scope}[xshift=7cm]
			% Draw the background
			\fill[goban] (-0.5,-0.5) rectangle (2.5,2.5);
			% Draw the grid lines
			\draw[step=1cm, black] (0,0) grid (2,2);
			% Draw the thick border
			\draw[very thick] (0,0) rectangle (2,2);
			\filldraw[draw=black,fill=white] (0,2) circle (12pt);
			\filldraw[draw=black,fill=white] (1,1) circle (12pt);
			\filldraw[black] (2,1) circle (12pt);
			\filldraw[black] (1,0) circle (12pt);
			\filldraw[black] (0,1) circle (12pt);
			\filldraw[draw=black,fill=white] (2,2) circle (12pt);
			\path (2,2) node [black]{1};
			\path (0,1) node [white]{2};
		\end{scope}
		\begin{scope}[xshift=10.5cm]
			% Draw the background
			\fill[goban] (-0.5,-0.5) rectangle (2.5,2.5);
			% Draw the grid lines
			\draw[step=1cm, black] (0,0) grid (2,2);
			% Draw the thick border
			\draw[very thick] (0,0) rectangle (2,2);
			\filldraw[draw=black,fill=white] (0,2) circle (12pt);
			\filldraw[draw=black,fill=white] (1,1) circle (12pt);
			\filldraw[black] (1,0) circle (12pt);
			\filldraw[black] (0,1) circle (12pt);
			\filldraw[draw=black,fill=white] (2,2) circle (12pt);
			\filldraw[draw=black,fill=white] (2,0) circle (12pt);
			\path (2,2) node [black]{1};
			\path (0,1) node [white]{2};
			\path (2,0) node [black]{3};
		\end{scope}
		\begin{scope}[xshift=14cm]
			% Draw the background
			\fill[goban] (-0.5,-0.5) rectangle (2.5,2.5);
			% Draw the grid lines
			\draw[step=1cm, black] (0,0) grid (2,2);
			% Draw the thick border
			\draw[very thick] (0,0) rectangle (2,2);
			\filldraw[draw=black,fill=white] (1,1) circle (12pt);
			\filldraw[black] (1,0) circle (12pt);
			\filldraw[black] (0,1) circle (12pt);
			\filldraw[draw=black,fill=white] (2,2) circle (12pt);
			\filldraw[draw=black,fill=white] (2,0) circle (12pt);
			\filldraw[black] (1,2) circle (12pt);
			\path (2,2) node [black]{1};
			\path (0,1) node [white]{2};
			\path (2,0) node [black]{3};
			\path (1,2) node [white]{4};
		\end{scope}
		\begin{scope}[xshift=0cm,yshift=-3.5cm]
			% Draw the background
			\fill[goban] (-0.5,-0.5) rectangle (2.5,2.5);
			% Draw the grid lines
			\draw[step=1cm, black] (0,0) grid (2,2);
			% Draw the thick border
			\draw[very thick] (0,0) rectangle (2,2);
			\filldraw[draw=black,fill=white] (1,1) circle (12pt);
			\filldraw[black] (0,1) circle (12pt);
			\filldraw[draw=black,fill=white] (2,2) circle (12pt);
			\filldraw[draw=black,fill=white] (2,0) circle (12pt);
			\filldraw[black] (1,2) circle (12pt);
			\filldraw[draw=black,fill=white] (0,0) circle (12pt);
			\path (2,2) node [black]{1};
			\path (0,1) node [white]{2};
			\path (2,0) node [black]{3};
			\path (1,2) node [white]{4};
			\path (0,0) node [black]{5};
		\end{scope}
		\begin{scope}[xshift=3.5cm,yshift=-3.5cm]
			% Draw the background
			\fill[goban] (-0.5,-0.5) rectangle (2.5,2.5);
			% Draw the grid lines
			\draw[step=1cm, black] (0,0) grid (2,2);
			% Draw the thick border
			\draw[very thick] (0,0) rectangle (2,2);
			\filldraw[draw=black,fill=white] (1,1) circle (12pt);
			\filldraw[black] (0,1) circle (12pt);
			\filldraw[draw=black,fill=white] (2,0) circle (12pt);
			\filldraw[black] (1,2) circle (12pt);
			\filldraw[draw=black,fill=white] (0,0) circle (12pt);
			\filldraw[black] (2,1) circle (12pt);
			\path (0,1) node [white]{2};
			\path (2,0) node [black]{3};
			\path (1,2) node [white]{4};
			\path (0,0) node [black]{5};
			\path (2,1) node [white]{6};
		\end{scope}
		\begin{scope}[xshift=7cm,yshift=-3.5cm]
			% Draw the background
			\fill[goban] (-0.5,-0.5) rectangle (2.5,2.5);
			% Draw the grid lines
			\draw[step=1cm, black] (0,0) grid (2,2);
			% Draw the thick border
			\draw[very thick] (0,0) rectangle (2,2);
			\filldraw[draw=black,fill=white] (1,1) circle (12pt);
			\filldraw[draw=black,fill=white] (2,0) circle (12pt);
			\filldraw[black] (1,2) circle (12pt);
			\filldraw[draw=black,fill=white] (0,0) circle (12pt);
			\filldraw[black] (2,1) circle (12pt);
			\filldraw[draw=black,fill=white] (0,2) circle (12pt);
			\path (2,0) node [black]{3};
			\path (1,2) node [white]{4};
			\path (0,0) node [black]{5};
			\path (2,1) node [white]{6};
			\path (0,2) node [black]{7};
		\end{scope}
		\begin{scope}[xshift=10.5cm,yshift=-3.5cm]
			% Draw the background
			\fill[goban] (-0.5,-0.5) rectangle (2.5,2.5);
			% Draw the grid lines
			\draw[step=1cm, black] (0,0) grid (2,2);
			% Draw the thick border
			\draw[very thick] (0,0) rectangle (2,2);
			\filldraw[draw=black,fill=white] (1,1) circle (12pt);
			\filldraw[black] (1,2) circle (12pt);
			\filldraw[draw=black,fill=white] (0,0) circle (12pt);
			\filldraw[black] (2,1) circle (12pt);
			\filldraw[black] (1,0) circle (12pt);
			\filldraw[draw=black,fill=white] (0,2) circle (12pt);
			\path (1,2) node [white]{4};
			\path (0,0) node [black]{5};
			\path (2,1) node [white]{6};
			\path (0,2) node [black]{7};
			\path (1,0) node [white]{8};
		\end{scope}
		\begin{scope}[xshift=14cm,yshift=-3.5cm]
			\path (1,1) node [black]{\Huge{\dots}};
		\end{scope}
	\end{tikzpicture}
	\caption{Ciclo de longitud 8 que se puede repetir, en ausencia de alguna regla de SuperKo.}
	\label{ciclo_inf}
\end{figure}
Podría argumentarse que, en el tablero de la Figura~\ref{ciclo_ko}, se ha repetido la posición, sí, pero el turno del jugador es distinto; por lo tanto, no se trata de una repetición completa. Este argumento conduce a la denominada regla del \textit{SuperKo Natural}. Sin embargo, nosotros haremos uso de una variante más restrictiva.
\begin{definition}[SuperKo Posicional]\label{sec-superko}
	Un jugador no puede colocar una piedra de manera tal que recree una posición que ya haya aparecido anteriormente en la partida.
\end{definition}

Por esta regla, tanto en el escenario de la Figura~\ref{ciclo_ko} como en el de la Figura~\ref{ciclo_inf}, cualquier intento de repetir un ciclo requerirá de jugadas intermedias que generen un avance real en el juego. Resta considerar la restricción más elemental del juego. No es posible jugar en un punto que ya se encuentre ocupado por otra piedra.

\begin{definition}[Condición de estado libre]
	Si se pretende registrar la jugada $f(n) = i$ en el turno $n$, es indispensable que el punto correspondiente se encuentre previamente libre. Esto se expresa como:
	\[f(n) = i \text{ es posicionable} \iff C(i) = 0.\]
\end{definition}
Ya tenemos el \textit{dónde}, \textit{quiénes} y \textit{qué}, veamos el \textit{cómo}. 

\subsection{Teorema de la jugada legal}

Hemos visto que una piedra $i$ es posicionable si $C(i)=0$. Una jugada no es suicida si $L(i) \neq \varnothing$. Para detectar una captura, debemos simular transitoriamente la jugada en $i$ para verificar si existe alguna piedra o grupo enemigo adyacente $X \in K(i)$ cuyo grupo se quede sin libertades. Esto sucede si $L(X)=\varnothing$, como la única libertad que le quedaba a ese grupo era, justamente, el punto donde acabamos de jugar, al ocuparla se concreta la captura. 

De haberse cumplido $L(X)=\varnothing$ antes del turno, el grupo ya habría sido retirado.
\[f(n) = i \text{ da lugar a una captura} \iff \exists X \in K(i) : L(X)=\varnothing.\]
La detección del \textit{SuperKo} requiere de un análisis considerablemente mayor, por lo que nos ocuparemos de ella en la Sección~\ref{detec_superko}.

Definimos las funciones booleanas \text{Libre}, \text{Suicidio}, \text{Captura} y \text{SKo}. Estas tomarán como argumento una potencial jugada $f(n) = i$ y devolverán $1$ si la condición se cumple, o $0$ en caso contrario:
\begin{align*}
	\text{Libre}(f(n)) &= \begin{cases}
		1 & \text{si } C(i) = 0 \\
		0 & \text{si } C(i) \neq 0.
	\end{cases} & \quad
	\text{Suicidio}(f(n)) &= \begin{cases}
		1 & \text{si } L(i) = \varnothing \\
		0 & \text{si } L(i) \neq \varnothing.
	\end{cases} \\[1em]
	\text{Captura}(f(n)) &= \begin{cases}
		1 & \text{si genera captura} \\
		0 & \text{si no}.
	\end{cases} & \quad
	\text{SKo}(f(n)) &= \begin{cases}
		1 & \text{si viola la regla} \\
		0 & \text{si no}.
	\end{cases}
\end{align*}

Con ellas podemos dar cuenta cuando una jugada es legal.

\begin{theorem}[Teorema de la jugada legal]\label{teoremalegal}
	Sea $f(n) = i$ una jugada en el turno $n$ de una partida de Go. Dicha jugada es legal si y solo si satisface la siguiente proposición lógica:
	\[ f(n) \text{ es legal} \iff \text{Libre} \wedge \neg \text{SKo} \wedge (\neg \text{Suicidio} \vee \text{Captura}). \]
\end{theorem}
\begin{proof}
	$(\Rightarrow)$ Supongamos que $f(n)$ es una jugada legal.
	Por definición, una jugada legal no viola la regla del \textit{SuperKo}, luego $\text{SKo}(f(n))=0$; por lo tanto, $\neg\text{SKo}(f(n))=1$.
	
	Además, sabemos que la \textit{regla del suicidio} prohíbe las jugadas suicidas a menos que estas produzcan una captura. Es decir, o bien la jugada no es un suicidio, o bien (si lo es) necesariamente produce captura. En términos lógicos, al ser $f(n)$ legal, la proposición $\neg\text{Suicidio} \vee \text{Captura}$ es siempre verdadera.
	
	La jugada debe realizarse sobre un punto vacío del tablero, es decir, $\text{Libre}(f(n))=1$.
	
	Solo existen tres tipos de jugadas legales:
	\begin{enumerate}
		\item El punto está libre, no hay \textit{SuperKo}, no hay suicidio y no se captura.
		\item El punto está libre, no hay \textit{SuperKo}, no hay suicidio y hay captura.
		\item El punto está libre, no hay \textit{SuperKo}, hay suicidio pero hay captura.
	\end{enumerate}
	
	$(\Leftarrow)$ Analicemos la tabla de verdad de nuestra proposición. Denotaremos \text{Suicidio} como $S$ y \text{Captura} como $\text{Cap}$ por cuestiones de legibilidad. Omitiremos los casos donde se viola el \textit{SuperKo} ($\text{SKo}=1$) o donde el punto está ocupado ($C(i) \neq 0$), dado que son trivialmente falsos.
	
	\[
	\begin{array}{c|c|c|c|c|c}
		\text{Libre} & \neg \text{SKo} & \neg S & \text{Cap} & \neg S \vee \text{Cap} & \text{Libre} \wedge \neg \text{SKo} \wedge (\neg S \vee \text{Cap}) \\ 
		\hline 
		1 & 1 & 1 & 1 & 1 & 1 \\
		1 & 1 & 1 & 0 & 1 & 1 \\
		1 & 1 & 0 & 1 & 1 & 1 \\
		1 & 1 & 0 & 0 & 0 & 0 \\
	\end{array}
	\]
	La tabla de verdad muestra que las únicas filas que hacen verdadera la proposición completa corresponden exactamente a las tres configuraciones deducidas en la ida.
\end{proof}
Más aún, no solo hemos caracterizado cuándo una jugada es legal, sino también qué tipos de jugadas existen:
\begin{enumerate}
	\item Colocar una piedra sin capturar.
	\item Realizar una captura.
	\item Realizar una captura mediante un suicidio.
\end{enumerate}

%		\subsection{SuperKo en dos dimensiones}
%		\label{SuperKo2d}
%		
%		Durante una partida, la detección visual de un \textit{ko} inmediato es natural; sin embargo, cuando enfrentamos una extensa seguidilla de jugadas que desemboca en una repetición cíclica del tablero, identificar la infracción puede volverse sumamente complejo. Para poder detectar cuándo ocurre esta repetición, necesitamos un método algorítmico que nos permita revisar el historial de la partida y determinar si, para un turno $t$, existió algún turno anterior $t-k$ en el que la posición se haya repetido.
%		
%		Si nuestro tablero es de $n \times m$, podemos representar su estado fácilmente mediante la siguiente matriz:
%		\[
%		\text{Go}(t) = \begin{bmatrix}
	%			e_{1,1} & e_{1,2} & \dots  & e_{1,m} \\
	%			e_{2,1} & e_{2,2} & \dots  & e_{2,m} \\
	%			\vdots  & \vdots  & \ddots & \vdots  \\
	%			e_{n,1} & e_{n,2} & \dots  & e_{n,m}
	%		\end{bmatrix}.
%		\]
%		Aquí, $\text{Go}(t)$ representa la matriz de estado de cada punto del tablero en el turno $t$.
%		
%		Supongamos que, en un turno $t$, nuestro tablero posee una piedra blanca en la posición $(1,1)$, adoptando la siguiente forma:
%		\[
%		\mathrm{Go}(t) = \begin{bmatrix}
	%			2 & e_{1,2} & \dots  & e_{1,m} \\
	%			e_{2,1} & e_{2,2} & \dots  & e_{2,m} \\
	%			\vdots  & \vdots  & \ddots & \vdots  \\
	%			e_{n,1} & e_{n,2} & \dots  & e_{n,m}
	%		\end{bmatrix}, \quad
%		\mathrm{Go}(t+k) = \begin{bmatrix}
	%			0 & e_{1,2} & \dots  & e_{1,m} \\
	%			e_{2,1} & e_{2,2} & \dots  & e_{2,m} \\
	%			\vdots  & \vdots  & \ddots & \vdots  \\
	%			e_{n,1} & e_{n,2} & \dots  & e_{n,m}
	%		\end{bmatrix}.
%		\]
%		Supongamos también que, luego de $k$ turnos, nos encontramos con un tablero $\mathrm{Go}(t+k)$ que es idéntico a $\mathrm{Go}(t)$, a excepción de esa primera posición, donde ahora $e_{1,1}(t+k) = 0$. Si en este momento es el turno del jugador blanco, sería ilegal que jugara en $(1,1)$, puesto que dicha acción viola la regla del \textit{SuperKo}.
\subsection{Contar los puntos}\label{contarpuntos}

En algún momento, por el SuperKo y mutuo acuerdo, la partida termina. Quien gana en el Go es el jugador que acumule la mayor cantidad de puntos. Para calcular esto, basta con sumar las piedras de cada jugador en el tablero al terminar la partida y sumar su territorio. Además, al jugador blanco se le suma el \textit{komi}.

Si definimos $P(2)$ como los puntos del blanco y $P(1)$ como los del negro, obtenemos:
\[P(1) = \sum_{I_1 \in I^{1}} |I_{1}| + \sum_{J_1 \in J^{1}} |J_{1}|, \quad P(2) = \sum_{I_2 \in I^{2}} |I_2| + \sum_{J_2 \in J^{2}} |J_{2}| + \text{komi}.\]

Para determinar al ganador, basta con evaluar la diferencia. Si $P(2) - P(1)\geq 0$, gana el blanco; en caso contrario, gana el negro.\footnote{Usamos $\geq$, pues contemplamos la posibilidad de que se decida jugar sin \textit{komi} o con un \textit{komi} de $0$ puntos, habilitando el empate.}

\section{El Go en \texorpdfstring{$n$}{n} dimensiones y \texorpdfstring{$m$}{m} jugadores}\label{Go_n_m}

Si pasamos a un espacio de $n$ dimensiones, necesitaremos $n$ conjuntos para poder construir el tablero. Tal como lo hicimos antes, para cada dimensión $i$, dado $n_i \in \mathbb{N}$, definimos $N_i = \{1, 2, \dots, n_i\}$. El tablero queda entonces conformado por el producto cartesiano:
\[ G = N_1 \times N_2 \times \dots \times N_n. \]
Para  $m$ jugadores, ampliamos el conjunto de estados a $E = \{0, 1, \dots, m\}$, reservando siempre el valor $0$ para las intersecciones libres.

\begin{definition}[Adyacencia entre puntos]
	Sea $T$ un tablero y, $t_1, t_2 \in T$ dos puntos dados por, $t_1 = (a_1, a_2, \dots, a_n),\ t_2 = (b_1, b_2, \dots, b_n).$
	
	Luego, $t_1$ es adyacente a $t_2$ y lo notaremos como $t_1 \sim t_2$ cuando,
	\[\exists! \ i \in \{1, 2, \dots, n\} : \big( \forall r \neq i, \ a_r = b_r \ \land \ a_i = b_i \pm 1 \big). \]
\end{definition}

Reducido a dos dimensiones, la condición exige una coordenada idéntica y otra difiriendo exactamente en 1, recuperando naturalmente la definición del tablero clásico.

De igual manera, extendemos la conexión entre piedras.
\begin{definition}[Conexión]
	Sea $T$ un tablero e $i_1, i_2 \in T$ tales que $i_1 = (e, t_1)$ e $i_2 = (e', t_2)$, estas están conectadas cuando:
	\[ i_1 \leftrightarrow i_2 \iff t_1 \sim t_2 \text{ y } C(i_1) = C(i_2). \]
\end{definition}
De ahora en adelante, cuando hagamos referencia a este tipo de conexión, la denominaremos \textit{adyacencia ortogonal}, puesto que veremos que es posible usar otros criterios en la Parte~\ref{parte3}.

Las definiciones de los grupos se mantienen idénticas, puesto que solo dependen de la conexión entre las piedras. A su vez, los territorios y las capturas funcionan exactamente igual, ya que solo necesitamos de las adyacencias del grupo. Lo único que realmente cambia al introducir más jugadores es el criterio de captura; más específicamente, la manera en que se distribuyen los puntos al retirar piedras enemigas. 

Antes de continuar, destaquemos una propiedad geométrica que es de vital importancia en las secciones siguientes.

\begin{lemma}\label{lem:2n}
	En un tablero de Go de $n$ dimensiones, un punto $t=(a_1,a_2,\dots,a_n)$ tiene al menos $n$ y a lo sumo $2n$ adyacencias.
\end{lemma}

\begin{proof}
	Para cada coordenada de $t$ existe al menos una dirección válida, pues si una está bloqueada por el borde la opuesta permanece disponible. En consecuencia cada dimensión aporta al menos una adyacencia. En un punto interior, cada coordenada aporta exactamente dos. Luego, hay al menos $n$ adyacencias y a lo sumo $2n$.
\end{proof}
Para simplificar la notación, convendremos en denotar al Go de $n$ dimensiones y $m$ jugadores como $\mathrm{Go}^n_m$. Si quisiéramos explicitar el tamaño de cada dimensión, podríamos hacerlo de la siguiente manera: $\mathrm{Go}^2_2(19,19)$ para el caso del tablero convencional. En caso de omitir la cantidad de jugadores, se interpretará que hablamos de un tablero de dimensión $n$ genérico, denotado por $\mathrm{Go}^n$. Del mismo modo, $\mathrm{Go}_m$ representará un tablero de dimensión arbitraria jugado por $m$ jugadores.
\subsection{Capturas y puntos para \texorpdfstring{$m$}{m} jugadores}\label{capturas_n_jug}

Sea $I$ un grupo de piedras. Así como en el tablero clásico, para que a dicho grupo se lo considere capturado, debe ocurrir que $L(I)=\varnothing$. Consideremos el tablero de $\mathrm{Go}^{2}_{4}$ ilustrado en la Figura~\ref{fig:capturasm}. Supongamos que nuestros jugadores son: celeste ($c$), rojo ($r$), verde ($v$) y marrón ($m$). Nuestro conjunto de estados queda definido entonces como $E = \{0, c, r, v, m\}$.

En este ejemplo, observamos que un grupo de 4 piedras rojas es capturado luego de que el jugador marrón coloca su piedra en la posición $(6,4)$. Sin embargo, notamos que el grupo rojo no fue rodeado en su totalidad de manera exclusiva por piedras marrones. De las 9 intersecciones adyacentes (las libertades que poseía originalmente el grupo rojo), 3 han sido ocupadas por el jugador verde, 2 por el celeste y, finalmente, 4 por el marrón.

¿Cómo distribuimos los puntos de captura en este escenario? Nuestra propuesta es repartir el puntaje de manera estrictamente proporcional a la influencia que cada color tuvo sobre el grupo capturado. Si el total a repartir son los 4 puntos correspondientes a las piedras rojas, los asignamos según la fracción de libertades ocupadas por cada atacante:
\begin{align*}
	\text{Puntos del verde} = \frac{3}{9} \cdot 4,\quad
	\text{Puntos del celeste} = \frac{2}{9} \cdot 4,\quad
	\text{Puntos del marrón} = \frac{4}{9} \cdot 4.
\end{align*}
En efecto, la suma de las partes es igual a la totalidad de los puntos:
\[ \left(\frac{3}{9} + \frac{2}{9} + \frac{4}{9}\right) \cdot 4 = 4. \]
\begin{figure}
	\centering
	\begin{tikzpicture}[scale=0.485] \tikzstyle{every node}=[font=\small]
		% Draw the background
		\fill[goban] (-0.5,-0.5) rectangle (6.5,6.5);
		
		% Draw the grid lines
		\draw[step=1cm, black] (0,0) grid (6,6);
		
		% Draw the thick border
		\draw[very thick] (0,0) rectangle (6,6);
		
		% Draw the central black dot
		
		\filldraw[Red] (2,4) circle (12pt);
		\filldraw[Red] (3,4) circle (12pt);
		\filldraw[Red] (4,3) circle (12pt);
		\filldraw[Red] (4,4) circle (12pt);
		\fill[RawSienna] (3,3) circle (12pt);
		\fill[TealBlue] (1,4) circle (12pt);
		\fill[TealBlue] (5,4) circle (12pt);
		\fill[OliveGreen] (3,5) circle (12pt);
		\fill[OliveGreen] (2,5) circle (12pt);
		\fill[OliveGreen] (4,5) circle (12pt);
		\fill[RawSienna] (4,2) circle (12pt);
		\fill[RawSienna] (2,3) circle (12pt);
		
		\begin{scope}[xshift=8cm]
			\fill[goban] (-0.5,-0.5) rectangle (6.5,6.5);
			
			% Draw the grid lines
			\draw[step=1cm, black] (0,0) grid (6,6);
			
			% Draw the thick border
			\draw[very thick] (0,0) rectangle (6,6);
			
			% Draw the central black dot
			
			%\filldraw[Red] (2,4) circle (12pt);
			%\filldraw[Red] (3,4) circle (12pt);
			%\filldraw[Red] (4,3) circle (12pt);
			%\filldraw[Red] (4,4) circle (12pt);
			\fill[RawSienna] (3,3) circle (12pt);
			\fill[TealBlue] (1,4) circle (12pt);
			\fill[TealBlue] (5,4) circle (12pt);
			\fill[OliveGreen] (3,5) circle (12pt);
			\fill[OliveGreen] (2,5) circle (12pt);
			\fill[OliveGreen] (4,5) circle (12pt);
			\fill[RawSienna] (4,2) circle (12pt);
			\fill[RawSienna] (2,3) circle (12pt);
			\fill[RawSienna] (5,3) circle (12pt);
			\path (5,3) node [white]{1};
		\end{scope}
	\end{tikzpicture}
	\caption{Captura en el $Go^{2}_{4}$. Juegan marrones y capturan al grupo rojo.}
	\label{fig:capturasm}
\end{figure}

De esta forma, se produce una distribución proporcional que recompensa a quien más contribuye a la captura. Otorgar los puntos solo a quien captura desincentivaría la agresión; sería inviable rodear a un oponente si en los $m-1$ turnos intermedios otros pudieran «robar» la captura.

\subsection{Komi para muchos jugadores}\label{komi_n_jug}

El origen del \textit{komi} radica en evitar empates y contrarrestar la desventaja de jugar segundo, por ello, es necesario reinterpretar este concepto cuando intervienen múltiples jugadores. Lo que sigue es una propuesta de cómo podríamos generalizarlo.
\subsection*{Komi escalonado}
Podemos definir el \textit{komi} como una función que depende estrictamente del orden del turno. Si consideramos el conjunto de jugadores $E \setminus \{0\}$, podemos establecer una constante de compensación $\alpha \in \mathbb{Q}$.  De esta forma, la función de \textit{komi} $K: E \setminus \{0\} \to \mathbb{R}$ se define como:
\[K(j) = (j - 1) \cdot \alpha.\]
Así, el primer jugador recibe $0$ puntos de compensación, el segundo recibe $\alpha$, el tercero $2\alpha$, y el último jugador recibe $(m-1)\alpha$.
\subsection*{Komi irracional}
Si además buscamos preservar la segunda función histórica del \textit{komi} — evitar empates —, podemos proseguir de la siguiente manera. Para garantizar que no existan dos jugadores con el mismo resultado final, podemos asignar un componente irracional único al \textit{komi} de cada participante. Dado que tanto el conteo de piedras y territorios (enteros) como las capturas proporcionales (racionales) se mantienen estrictamente dentro de $\mathbb{Q}$, la adición de un valor irracional específico para cada jugador rompe cualquier posibilidad de igualdad. Al ser la suma de un racional y un irracional es siempre irracional, la igualdad entre dos puntajes es imposible pues sus irracionales son siempre distintos.

\subsection*{Komi mediante subasta}
Antes de iniciar la partida, cada jugador $j \in E \setminus \{0\}$ declara de forma secreta una apuesta $b_j \in \mathbb{Q}$, que representa la cantidad de puntos que está dispuesto a sacrificar de su recuento final a cambio del privilegio de asegurar la iniciativa, es decir un \textit{komi} negativo. Una vez se tienen todas las apuestas, se revelan. Se establece el orden de los turnos de manera estrictamente decreciente: el jugador con el valor $b_j$ más alto jugará primero, seguido por el segundo más alto, y así sucesivamente. En caso de igualdad entre dos o más apuestas, el orden de esos jugadores se resuelve a favor de quien apostó primero.

Estas soluciones pueden usarse en simultáneo o bien ser reforzadas, lo aquí presentado son ideas de cómo poder subsanar la carencia de un \textit{komi} bien definido en el Go generalizado.\footnote{Nota de autor: A mi parecer el \textit{komi} no debería existir, la iniciativa (\textit{sente}) puede perderse en cualquier momento. Además, mientras mayor es el tablero, menos importa quien juega primero. Esta invención del siglo XX solo es útil para jugadores profesionales.}

\section{Detección del SuperKo posicional}\label{detec_superko}

Gracias al Teorema de la jugada legal~\ref{teoremalegal}, el cual por su construcción vale para $n$ dimensiones, poseemos la herramienta para saber si una jugada es legal o no; sin embargo, hasta ahora no hemos discutido el mecanismo para detectarlo.

Ya hemos formalizado cómo determinar si hay captura, si hay suicidio, o si el punto deseado se encuentra ocupado. Ahora bien, ¿Cuándo hay \textit{SuperKo}? Es decir, ¿cómo sabemos de manera efectiva que hemos roto esta regla?

Que exista una violación a la \textit{regla del SuperKo} implica que la jugada propuesta recrea una configuración del tablero que ya ha ocurrido en la partida. Para detectar esto, bastaría con poseer un registro histórico de los estados de cada punto en cada turno. De esta manera, simulando el turno que se pretende jugar y comparándolo con el historial, sabríamos que, si existe una coincidencia exacta, se ha roto la regla.

Para resolver este problema, haremos uso del trabajo desarrollado por Albert L. Zobrist \cite{zobrist1970hashing}. 

%		El primer método propone lo siguiente: dado un tablero $\mathrm{Go}^{n}_{m}$, buscamos simplificar la representación espacial de las $n$-tuplas y colapsar todas las coordenadas en un vector. El tamaño o dimensión de este vector será exactamente el cardinal del tablero:
%		\[|T| = \prod_{i=1}^{n} |N_i| = \tau.\]
%		Basta extraer el estado de cada intersección para obtener una $\tau$-tupla que codifica el tablero de manera \textit{ordenada}. Antes de ejecutar una jugada, se simula el tablero resultante y se realiza la siguiente resta con el historial, si convenimos en llamara a este historial $T(t)$, el tablero en el turno $t$, queremos comparar lo siguiente,
%		\[T(t) - T(t-k).\]
%		Para $k = 1, 2, \dots, t-1$. Si el resultado de esta operación es el vector nulo para algún valor de $k$, entonces la posición es idéntica y hemos roto la regla del \textit{SuperKo}. 


%		Veremos que este proceso puede optimizarse considerablemente en términos de la cantidad de comparaciones requeridas.
%		
%		\subsection{Función emparejadora de Cantor}\label{functCantor}
%		
%		Para estos efectos, haremos uso de la función emparejadora de Cantor generalizada (\textit{cpf}).
%		
%		\begin{theorem}
	%			Para todo $n \in \mathbb{N}$, la función emparejadora de $n$-tuplas $\langle \cdot, \dots, \cdot \rangle: \mathbb{N}^{n} \rightarrow \mathbb{N}$, definida como:
	%		\[cpf(X_1, \dots, X_n) = \sum_{h=1}^{n} \left\{ \frac{1}{h!} \prod_{j=0}^{h-1} \left[ \left( \sum_{i=1}^{h} X_i \right) + j \right] \right\}, \quad \forall(X_{1}, \dots, X_{n}) \in \mathbb{N}^{n}.\]
	%			es un polinomio de grado $n$ con coeficientes racionales. Más aún, es una biyección de $\mathbb{N}^n \rightarrow \mathbb{N}$.\footnote{Para ver la demostración, diríjase a \cite{Lisi2007}.}
	%		\end{theorem}
%		
%		Por el teorema, todo tablero puede ser llevado a $\mathbb{N}$. Al ser nuestros tableros finitos, aplicar el teorema produce en una inyección de $G$ en $\mathbb{N}$. Si nuestro tablero $T$ es un subconjunto de $E \times G$, con $E = \{0, 1, 2, \dots, m\}$ y $G = N_1 \times \dots \times N_n$, definiendo el número total de intersecciones como $\pi = |G| = \prod_{i=1}^{n} |N_i|$, obtenemos la siguiente correspondencia:
%		\[ T \leftrightarrow E \times \mathbb{N}_{\pi} = \{(e_1, k_1), (e_2, k_2), \dots, (e_{\pi}, k_\pi)\} \]
%		donde $e_i \in E$ es algún estado y $k_j$ es un número natural único dado por la función. Cabe aclarar que aquí $\mathbb{N}_{\pi}$ no representa el conjunto de los naturales «hasta el número $\pi$», sino un subconjunto específico de números naturales de cardinalidad $\pi$.
%		
%		Con esto, logramos ordenar el tablero: basta con seguir el orden natural de $\mathbb{N}$ en la segunda coordenada. Como lo único que nos importa para la detección de la regla son los estados del tablero (y no los números en sí, podemos dar un paso más. Si aplicamos la función de estado $C$ que definimos en la Sección~\ref{funcioncolor}, obtenemos un vector de dimensión $\pi$:
%		\[ C \circ \textit{cpf} (E \times \mathbb{N}_{\pi}) = (e_{1}, e_{2}, \dots, e_{\pi}) = E^{\pi}. \]
%		Valiéndonos de la biyección, este vector está perfectamente definido, pues a cada piedra $i=(e,t)$ le corresponde un único par $(e, \textit{cpf}(t))$. Esta correspondencia se ilustra en el diagrama conmutativo de la Figura~\ref{diagramacantor}.
%		
%		\begin{figure}[H]
	%			\centering
	%			\[
	%			\begin{tikzcd}
		%				{E \times G} && {E \times \mathbb{N}_{\pi}} \\
		%				\\
		%				{E^{\pi}}
		%				\arrow["\textit{cpf}", leftrightarrow, from=1-1, to=1-3]
		%				\arrow["C"', from=1-1, to=3-1]
		%				\arrow["{C \circ \textit{cpf}}", from=1-3, to=3-1]
		%			\end{tikzcd}
	%			\]
	%			\caption{Diagrama de la composición.}
	%			\label{diagramacantor}
	%		\end{figure}
%		
%		Como el tablero posee un ordenamiento estricto y la función es biyectiva, dos tableros idénticos siempre generarán el mismo vector $E^{\pi}$.
%		
%		Si guardamos el vector $E^{\pi}(t)$ en cada turno $t$, para verificar la legalidad de la jugada actual basta restar $E^{\pi}(t) - E^{\pi}(t-k)$ para todo $k = 1, 2, \dots, t-1$. Si para algún $k$ dicha resta da como resultado el vector nulo ($\vec{0}$), significa que la posición exacta se ha repetido, rompiendo la regla. Formalmente:
%		\[ \text{SuperKo en el turno } t \iff \exists k : E^{\pi}(t) - E^{\pi}(t-k) = \vec{0}. \]
%		Como hemos adelantado, no hace falta realizar las $t-1$ comparaciones. Para que un tablero pueda siquiera tener la posibilidad de ser igual a otro, es una condición necesaria que la cantidad de piedras de cada color sea idéntica en ambos turnos. Es decir, solo computaremos $E^{\pi}(t) - E^{\pi}(t-k)$ si se cumple que, para cada estado $e \in E \setminus \{0\}$:
%		\[ \sum_{I} |I^{e}|_t = \sum_{I} |I^{e}|_{t-k}. \]
%		
%		Podría parecer que estamos realizando cálculos de más, pero el tamaño del tablero ($\pi$) siempre será inmensamente mayor que la cantidad de jugadores ($m$). Guardar el conteo de las piedras exige almacenar solo $m$ números adicionales por turno. Realizar a lo sumo $m$ comparaciones numéricas simples nos permite descartar rápidamente un tablero que, de otro modo, exigiría $\pi$ comparaciones componente a componente. A modo de ejemplo, en el $\mathrm{Go}^{2}_{2}(19,19)$, basta hacer 2 comparaciones preliminares (conteo de blancas y negras) para descartar el análisis de un tablero de $19 \cdot 19 = 361$ puntos. Por lo tanto, nos ahorramos hasta 359 operaciones por turno descartado.
%Veamos un ejemplo con el tablero tradicional. Si tenemos $E\times G$, $E=\{a,e_1,e_2\}$, $G=N\times N$ con $|N|=3$, la función de cantor para 2 dimensiones nos queda:

%$$\langle a,b \rangle =  \sum_{h=1}^{2}\left\lbrace \frac{1}{h!} \prod_{j=0}^{h-1} \left[\ \left( \sum_{i=1}^{h} X_i \right)+j \right]  \right\rbrace=\frac{(a+b)²+3a+b}{2}$$

%Con lo cual, obtenemos la siguiente tabla que corresponde a los puntos evaluados en la función.

%\begin{table}
%	\centering
%		\begin{tabular}{|l|l|ll}
	%			\hline
	%			cdf(x,y) & Número & \multicolumn{1}{l|}{cdf(x,y)}&
	%\multicolumn{1}{l|}{Número} \\ \hline
	%			1,1   & 4      & \multicolumn{1}{l|}{2,3}      & \multicolumn{1}{l|}{18}     \\ \hline
	%			1,2      & 8      & \multicolumn{1}{l|}{3,1}      & \multicolumn{1}{l|}{11}     \\ \hline
	%			1,3      & 13     & \multicolumn{1}{l|}{3,2}      & \multicolumn{1}{l|}{17}     \\ \hline
	%			2,1      & 7      & \multicolumn{1}{l|}{3,3}      & \multicolumn{1}{l|}{24}     \\ \hline
	%			2,2      & 12     &                               &                             \\ \cline{1-2}
	%		\end{tabular}
%	\end{table}

% Obtenemos para cada turno el conjunto (ordenado)
%\begin{align*}
%	E\times M(t)=\{(e_{1},4),(e_{2},7),(e_{3},8),(e_{4},11),(e_{5},12)\\,(e_{6},13),(e_{7},17),(e_{8},18),(e_{9},24)\}.
%\end{align*}

%Con $e_i\in E$, Que podemos simplificar a $E(t)=(e_{1},\dots,e_{9}).$ Entiéndase que los $e_i$ no están dados por el orden, sino que están numerados para mejor legibilidad.
% Que corresponde al tablero dado por

% \begin{align*}
%	E\times G=\{(e_{i_{1}},(1,1)),(e_{i_{2}},(2,1)),(e_{i_{3}},(1,2)),(e_{i_{4}},(3,1)),\\
%	(e_{i_{5}},(2,2)),(e_{i_{6}},(1,3)),(e_{i_{7}},(3,2)),(e_{i_{8}},(2,3)),(e_{i_{9}},(3,3))\}
%\end{align*}

%	  Si bien utilizar \textit{cpf} nos proporciona una biyección y, por lo tanto, una representación infalible para nuestros fines, resulta computacionalmente muy costoso. Evaluar el polinomio para cada posición en tableros masivos y realizar múltiples comparaciones vectoriales en partidas largas es ineficiente.

%	  Afortunadamente, existe otra forma de hacer todo esto en tiempo constante ($O(1)$), con el único inconveniente de que renuncia a la biyección absoluta y adquiere una naturaleza probabilística (puede fallar por colisiones). 


\subsection{Hash de Zobrist}\label{ZobristHash}

Para este método, haremos uso del operador lógico \textit{o exclusivo}, denotado como XOR ($\oplus$), aplicado sobre números aleatorios $r_i$ representados como enteros de $n$ bits.

Un entero de $n$ bits es, una secuencia de $n$ dígitos binarios. Por ejemplo, el número 15 en 4 bits se representa como $1111$, puesto que:
\[ 2^3 \cdot 1 + 2^2 \cdot 1 + 2^1 \cdot 1 + 2^0 \cdot 1 = 8 + 4 + 2 + 1 = 15. \]
Del mismo modo, el número 9 en 4 bits es $1001$. Si recordamos la tabla de verdad para el operador lógico XOR, Figura~\ref{TablaXOR}.
\vspace{10pt}

\begin{wrapfigure}[9]{l}{0.43\textwidth}
	\vspace{-20pt}
	%\begin{figure}[H]
	\centering
	\[
	\begin{array}{c|c|c}
		A & B & A \oplus B \\
		\hline
		0 & 0 & 0 \\
		0 & 1 & 1 \\
		1 & 0 & 1 \\
		1 & 1 & 0
	\end{array}
	\]
	\caption{Tabla de verdad del XOR.}
	\label{TablaXOR}
\end{wrapfigure}Realizar la operación XOR entre 15 y 9 en base 2 significa aplicar el operador de manera individual a cada bit correspondiente. Vale aclarar que, si bien escribimos el número de forma contigua, en realidad lo tratamos como si fueran, en este caso, 4-tuplas: $1111 \oplus 1001 = 0110$. Si $r_i$ y $r_j$ son dos secuencias aleatorias de enteros de $n$ bits, asumiremos que se cumplen las siguientes propiedades:\hfill\break

\begin{enumerate}
	\item[P1)] $( r_i \oplus r_j ) \oplus r_k = r_i \oplus ( r_j \oplus r_k )$. \hfill(Asociatividad)
	\item[P2)] $r_i \oplus r_j = r_j \oplus r_i$.\hfill (Conmutatividad)
	\item[P3)] $r_i \oplus r_i = 0$.\hfill (nilpotencia)
	\item[P4)] Si $s = r_1 \oplus r_2 \oplus \dots \oplus r_k$, entonces $s$ es también una secuencia aleatoria.
	\item[P5)] Los valores resultantes $s$ se encuentran uniformemente distribuidos.
\end{enumerate}

Para comenzar a construir nuestro algoritmo, debemos tomar el conjunto de todas las configuraciones posibles de una piedra en el tablero, es decir, $E \setminus \{0\} \times G.$
Donde excluimos el $0$ puesto que no colocamos piedras de color $0$ al jugar. Luego, asignamos un número aleatorio de $n$ bits, único, a cada elemento. Si recordamos que un jugador posee la opción de \textit{pasar} su turno (una jugada que no altera el estado del tablero), a esta acción se le asignará $\vec{0}$, lo que garantiza que el estado no cambie.

Para inicializar el hash, debemos garantizar más números que posiciones. Conociendo la cardinalidad $|E\setminus\{0\}\times G|$ y sumando 1 para el tablero vacío, la longitud mínima de n es:
\[ 2^n \geq |E \setminus \{0\} \times G| + 1. \]
Para nuestro caso, en el tablero convencional, tendremos $361$ intersecciones y $2$ estados posibles para las piedras (blanco y negro). Por lo tanto:
$$|E\setminus\{0\} \times G| = 2 \times 361 = 722.$$
Luego, el número mínimo de bits necesario para mapear el tablero base será $n = 10$, puesto que
$$2^{10} = 1024 \geq 722 + 1.$$

Zobrist advierte sobre la existencia de dos tipos de errores. Un \textit{error Tipo 1} ocurre cuando dos configuraciones distintas del tablero producen exactamente el mismo valor de hash; en el contexto del \textit{SuperKo}, esto resultaría en un falso positivo. Por otro lado, un \textit{error Tipo 2} se produce debido a limitaciones de espacio en la memoria, lo cual no consideraremos.

Para reducir la posibilidad de un error Tipo 1, Zobrist propone dos soluciones complementarias. La primera consiste en generar un segundo hash en simultáneo, utilizando una secuencia independiente, para descartar los falsos positivos. La segunda solución radica en utilizar una mayor cantidad de bits de los que parecen necesarios. Si $n$ es la cantidad de bits que pretendíamos usar y $m$ los bits a agregar, Zobrist demuestra que la probabilidad de que no ocurran errores de Tipo 1 viene dada por la siguiente expresión (Tabla~\ref{Tabla_prob}):
$$p(\text{no error}) = \left( 1 - \frac{1}{2^{n+m}} \right)^{2^n} \approx \left( \frac{1}{e} \right)^{2^{-m}}.$$
\begin{table}[H]
	\centering
	\begin{tabular}{c|c || c|c}
		\textbf{$m$ (bits extra)} & \textbf{Prob. de no error} & \textbf{$m$ (bits extra)} & \textbf{Prob. de no error} \\ \hline
		0 & 0.36778 & 6  & 0.98449 \\
		1 & 0.60653 & 7  & 0.99222 \\
		2 & 0.77879 & 8  & 0.99611 \\
		3 & 0.88251 & 9  & 0.99805 \\
		4 & 0.93941 & 10 & 0.99901 \\
		5 & 0.96924 & \multicolumn{2}{c}{}
	\end{tabular}
	\caption{Probabilidad de no error al utilizar más bits.}
	\label{Tabla_prob}
\end{table}
Por lo tanto, bastaría con utilizar números de 21 o incluso 32 bits para lograr que la probabilidad de evitar un error de Tipo 1 sea de 0.9999995 o 0.9999999997, respectivamente. Podríamos generar dos, tres o cuantas secuencias de verificación consideremos necesarias, con el fin de minimizar el error.

Para verificar la regla del \textit{SuperKo}, basta con realizar $S \cdot (t-1)$ operaciones, donde $S$ es el número total de secuencias hash independientes implementadas. Si $H_{s}(t)$ representa el valor del hash en el turno $t$ para la secuencia $s$, entonces:
$$ \text{SuperKo en el turno } t \iff H_{s}(t) \oplus H_{s}(t-k) = \vec{0}\ \ \text{y}\ \  f(t),f(t-k)\neq \varnothing.$$
para algún $k \in \{1, 2, \dots, t-1\}$ y para toda secuencia $s \in \{1, \dots, S\}$. Los números que asignamos a las piedras, los notaremos por por $h_s(i)$. Además, para procesar las capturas, haremos uso de las propiedades del operador XOR que asumimos previamente. Cuando una piedra es capturada, esta es retirada del tablero. Por lo tanto, si la piedra $p_{j}$ es removida, basta con aplicar la operación $H_s(t) \oplus h_s(i_{j})$ al estado actual.

Gracias a las Propiedades P1 y P2, podemos reordenar el cálculo del hash:
\begin{align*}
	H_s(t) \oplus h_s(i_{j}) &= \big( h_s(i_{1}) \oplus \dots \oplus h_s(i_{j}) \oplus \dots \oplus h_s(i_{m}) \big) \ \textcolor{Red}{\oplus \ h_s(i_{j})} \\
	H_s(t) \oplus h_s(i_{j}) &= h_s(i_{1}) \oplus \dots \oplus \big( h_s(i_{j}) \ \textcolor{Red}{\oplus \ h_s(i_{j})} \big) \oplus \dots \oplus h_s(i_{m}) & \text{(P1 y P2)} \\
	H_s(t) \oplus h_s(i_{j}) &= h_s(i_{1}) \oplus \dots \oplus \textcolor{Red}{\vec{0}} \oplus \dots \oplus h_s(i_{m}). & \text{(P3)}
\end{align*}

Dado que el vector nulo $\vec{0}$ es el elemento neutro de la operación $(\vec{0}\oplus h(i_{p})=h(i_{p}))$, el término desaparece. Así, demostramos que retirar una piedra (capturar) equivale a sumar nuevamente su hash. Para cerrar este capítulo, ilustremos con un ejemplo concreto en un tablero pequeño de $\mathrm{Go}^{2}_{2}(3,3)$.

Como el tablero posee 9 puntos, tenemos $|(E\setminus\{0\}) \times G| = 18$. Luego, el número mínimo de bits $n$ necesario será 5, puesto que $2^5 = 32 \geq 18 + 1$. Si decidimos agregar 3 bits adicionales como margen de seguridad, tendremos números de 8 bits. Para inicializar, basta con asignar de manera aleatoria un número a cada estado posible en cada punto. A tales efectos, construimos la siguiente tabla:
\begin{table}[H]
	\centering
	\small
	\begin{tabular}{lc|lc|lc}
		\textbf{Piedra} & \textbf{Número} & \textbf{Piedra} & \textbf{Número} & \textbf{Piedra} & \textbf{Número} \\ \hline
		(1,(1,1)) & 01111111 & (1,(2,1)) & 00010110 & (1,(3,1)) & 01001101 \\
		(1,(1,2)) & 10001101 & (1,(2,2)) & 11110001 & (1,(3,2)) & 11101111 \\
		(1,(1,3)) & 10111110 & (1,(2,3)) & 10110001 & (1,(3,3)) & 00011110 \\
		(2,(1,1)) & 00001110 & (2,(2,1)) & 00000010 & (2,(3,1)) & 01001101 \\
		(2,(1,2)) & 10010011 & (2,(2,2)) & 10111010 & (2,(3,2)) & 01000111 \\
		(2,(1,3)) & 10110111 & (2,(2,3)) & 10010001 & (2,(3,3)) & 11001100 \\
	\end{tabular}
	\caption{Tabla de los hashes de cada punto del tablero de $3\times 3$}
	\label{tablahash}
\end{table}

Inicializamos el hash de la partida en el vector nulo $\vec{0} = 00000000$, representado el tablero vacío. A partir de aquí, por cada jugada efectuada, actualizamos el estado aplicando la operación XOR ($\oplus$) entre el hash del turno actual y el valor de la nueva piedra colocada (o retirada, en caso de captura). Supongamos que se realiza la siguiente secuencia de jugadas,
obteniendo el tablero de la Figura~\ref{fig:tablero_ejsko}: \begin{align*}
	f(1)&=(1,(2,2)), &f(2)&=(2,(1,3)),	&f(3)&=(1,(2,3)), &f(4)&=(2,(2,1)),\\
	f(5)&=(1,(3,1)), &f(6)&=(2,(3,2)), &f(7)&=(1,(1,2)), &f(8)&=\varnothing,\\
	f(9)&=\varnothing. 
\end{align*} 
%\begin{wrapfigure}[4]{r}{0.5\textwidth}
\begin{figure}[H]
	\centering
	\begin{tikzpicture}[scale=0.75]
		
		
		% Draw the background
		\fill[goban] (-0.5,-0.5) rectangle (2.5,2.5);
		% Draw the grid lines
		\draw[step=1cm, black] (0,0) grid (2,2);
		% Draw the thick border
		\draw[very thick] (0,0) rectangle (2,2);
		
		% Draw the central black dot
		\filldraw[draw=black,fill=white] (0,2) circle (12pt);
		\filldraw[draw=black,fill=white] (1,0) circle (12pt);
		%\filldraw[draw=black,fill=white] (2,1) circle (12pt);
		\filldraw[black] (1,1) circle (12pt);
		\filldraw[black] (1,2) circle (12pt);
		\filldraw[black] (2,0) circle (12pt);
		\path (1,1) node [white]{1};
		\path (0,2) node [black]{2};
		\path (1,2) node [white]{3};
		\path (1,0) node [black]{4};
		\path (2,0) node [white]{5};
		
		\begin{scope}[xshift=3.5cm]
			% Draw the background
			\fill[goban] (-0.5,-0.5) rectangle (2.5,2.5);
			% Draw the grid lines
			\draw[step=1cm, black] (0,0) grid (2,2);
			% Draw the thick border
			\draw[very thick] (0,0) rectangle (2,2);
			
			% Draw the central black dot
			\filldraw[draw=black,fill=white] (0,2) circle (12pt);
			\filldraw[draw=black,fill=white] (1,0) circle (12pt);
			\filldraw[draw=black,fill=white] (2,1) circle (12pt);
			\filldraw[black] (1,1) circle (12pt);
			\filldraw[black] (1,2) circle (12pt);
			\path (1,1) node [white]{1};
			\path (0,2) node [black]{2};
			\path (1,2) node [white]{3};
			\path (1,0) node [black]{4};
			\path (2,1) node [black]{6};
			
		\end{scope}
		\begin{scope}[xshift=7cm]
			% Draw the background
			\fill[goban] (-0.5,-0.5) rectangle (2.5,2.5);
			% Draw the grid lines
			\draw[step=1cm, black] (0,0) grid (2,2);
			% Draw the thick border
			\draw[very thick] (0,0) rectangle (2,2);
			
			% Draw the central black dot
			%\filldraw[draw=black,fill=white] (0,2) circle (12pt);
			\filldraw[draw=black,fill=white] (1,0) circle (12pt);
			\filldraw[draw=black,fill=white] (2,1) circle (12pt);
			\filldraw[black] (1,1) circle (12pt);
			\filldraw[black] (1,2) circle (12pt);
			\filldraw[black] (0,1) circle (12pt);
			\path (1,1) node [white]{1};
			\path (0,1) node [white]{7};
			\path (1,2) node [white]{3};
			\path (1,0) node [black]{4};
			\path (2,1) node [black]{6};
		\end{scope}
	\end{tikzpicture}
	\caption{Partida terminada en un tablero de $3\times 3$.}
	\label{fig:tablero_ejsko}
\end{figure}

Calcularemos los estados turno a turno, separando los vectores en bloques de 4 bits para facilitar su lectura. Si denotamos el hash de una piedra como $h(p_i)$, donde $p_i = (e, t)$, la evolución de la partida sería como sigue. En la sexta jugada se efectúa una captura sobre la piedra negra en $(3,1)$, por lo tanto, además de aplicar el hash de la piedra jugada, debemos aplicar el hash de la piedra retirada para «deshacer» su presencia. Si se capturaran múltiples piedras, bastaría con iterar la operación para cada una de ellas en cualquier orden (por P1, P2):
\begin{figure}[H]
	\centering
	\begin{align*}
		H(0) &= \vec{0} \\
		H(1) &= H(0) \oplus \textcolor{Blue}{h((1, (2,2)))} = 1111\, 0001 \\
		H(2) &= H(1) \oplus \textcolor{Blue}{h((2, (1,3)))} = 0100\, 1111 \\
		H(3) &= H(2) \oplus \textcolor{Blue}{h((1, (2,3)))}= 1011\, 1110 \\
		H(4) &= H(3) \oplus \textcolor{Blue}{h((2, (2,1)))} = 1010\, 1000 \\
		H(5) &= H(4) \oplus \textcolor{Blue}{h((1, (3,1))) }= 0110\, 0101 \\
		H(6) &= H(5) \oplus \textcolor{Blue}{h((2, (3,2)))} \oplus \textcolor{Red}{h((1, (3,1)))} = 1100\, 0111 \\
		H(7)  &= H(6) \oplus \textcolor{Blue}{h((1, (1,2)))} \oplus \textcolor{Red}{h((2, (1,3)))}= 1111\, 0100 \\
		H(8)  &= H(7) \oplus \textcolor{Blue}{0000\, 0000} = H(7) \\
		H(9)  &= H(8) \oplus  \textcolor{Blue}{0000\, 0000} = H(8).
	\end{align*}
	\caption{Hashes de la partida, en azul las \textcolor{Blue}{jugadas}, en rojo las \textcolor{Red}{capturas}.}
\end{figure} 
En los turnos 8 y 9, pasar el turno produce un estado idéntico al anterior. Esto no viola el \textit{SuperKo}, pues la regla prohíbe unicamente colocar una piedra para recrear una posición.

Ahora bien, podría suceder que, $h(2,(3,2))=1100\ 1011$. Provocando en el sexto turno un falso positivo para Superko, pues $H(6)= H(5)\oplus h(2,(3,2)) = 1011\ 1110.$ y comparando con el historial, $$H(6)\oplus H(3)=0000\ 0000.$$
De esta manera, con más bits e implementando múltiples secuencias simultáneas para anular las colisiones, el método de Zobrist queda implementado en el Go.

\section{La estructura del tiempo}

Ya estamos en condiciones de inyectar el tiempo en la estructura. Definimos a $\mathcal{T}$ como el espacio de todos los tableros,
$$ \mathcal{T} = \{ T \subset E \times G \mid T \text{ es un tablero}\}.$$

\subsection{La partida como trayectoria}
Una partida resulta una sucesión de jugadas $f:\mathbb{N} \rightarrow \mathcal{J}$, donde $f(t)$ representa la acción elegida por el jugador activo en el turno $t$. Esta función $f$ induce de manera natural una sucesión de tableros en $\mathcal{T}$, a la que denominaremos trayectoria:
$$\{T_0, T_1, \dots, T_{fin}\} $$

La sucesión se construye de manera recursiva a través de $\Delta$:
\begin{enumerate}
	\item \textbf{Estado inicial ($t=0$):} $T_0$ es el tablero vacío, es decir, el conjunto $T_0 = \{(0, t) : t \in G\}$.
	\item \textbf{Transición ($t \rightarrow t+1$):} El tablero en el turno siguiente es el resultado de aplicar la transición a la jugada del turno actual:
	$$ T_t = \Delta(T_{t-1}, f(t)) $$
\end{enumerate}

\subsection{Dinámica del estado}

Bajo esta construcción, formalizamos las propiedades del flujo del juego. Denotaremos como $E^k(T_t) = \{i \in T_t : C(i) = k\}$ al conjunto de todas las piedras de color $k$ presentes en el tablero en el turno $t$, y como $|E^k(T_t)|$ a su cantidad.

\begin{itemize}
	\item \textbf{Crecimiento y decrecimiento:} Si en el turno $t$ el jugador activo es $k \in E \setminus \{0\}$, el nuevo tablero dado por $T_t = \Delta(T_{t-1}, f(t))$, poseerá, a lo sumo, una piedra más de su color. En contraste, las piedras de los demás colores ($j \neq k$) solo pueden decrecer o mantenerse constantes debido a posibles capturas:
	\begin{align*}
		|E^k(T_t)| &\le |E^k(T_{t-1})| + 1,\\
		|E^j(T_t)| &\le |E^j(T_{t-1})|, \quad \forall j \neq k.
	\end{align*}
	
	
	\item \textbf{Condición de parada:} Sea $m = |E \setminus \{0\}|$ la cantidad total de jugadores. La partida concluye en el turno $T_{fin}$ cuando todos los jugadores pasan consecutivamente. Esto ocurre cuando:
	$$f(T_{fin}) = f(T_{fin-1}) = \dots = f(T_{fin-m+1}) = \varnothing$$
	En consecuencia, el estado del tablero se congela, resultando en:
	$$T_{{fin}} = T_{{fin-1}} = \dots = T_{{fin-m}}$$
	Este estado $T_{{fin}}$ constituye el Tablero Final.
\end{itemize}

\subsection{Anatomía de una jugada}
Entendamos cómo funciona $\Delta$. Supongamos que estamos jugando al Go y que, en nuestro turno, pretendemos jugar $f(t)$.
Si $f(t) = \varnothing$, el tablero permanece idéntico y pasamos el turno, es decir,
$$T_{t} = \Delta(T_{t-1}, \varnothing) = T_{t-1}.$$
Si $f(t) = i^{*} = (e, p)$, hemos de validar dicha jugada, el cálculo del nuevo tablero $T_t = \Delta(T_{t-1}, i^{*})$ en virtud del teorema~\ref{teoremalegal} se realiza de la siguiente manera:
\begin{enumerate}
	\item \textbf{Colocación:} Si existe $(e',p)\in T_{t-1}$ tal que $e' \neq 0$, la jugada es ilegal. Si $e' =0$, generamos un tablero transitorio $T'$ donde el punto $p$, previamente vacío, adquiere el estado $e$. Luego avanzamos al siguiente paso.
	$$ T' = \left(T_{t-1} \setminus \{(0, p)\}\right) \cup \{(e, p)\}.$$
	\item \textbf{Detección de SuperKo:}
	Calculamos $H(T_{t-1}) \oplus h(i^{*})$ y lo comparamos con el historial. Si su comparación da $\vec{0}$, en cada hash, hemos roto la regla y la jugada se descarta, en caso contrario avanzamos al siguiente paso.
	\item \textbf{Detección de captura:} Evaluando las adyacencias en $T'$, identificamos el conjunto $Cap$, uniendo todos los grupos, de adversarios, que no poseen libertades.
	$$ Cap = \bigcup \{ i\in I^{e'} \subset T' \mid L(I^{e'}) = \varnothing, e' \notin \{0, e\}\}.$$
	Si $Cap \neq \varnothing$, la jugada es legal y saltamos al paso $5$, en caso contrario avanzamos al siguiente paso.
	\item \textbf{Detección de suicidio:} Considerando $Cap = \varnothing$, evaluamos el estado de la piedra jugada $i^* = (e, p)$ en $T'$. El análisis se divide en dos casos:
	\begin{itemize}
		\item \textbf{Piedra aislada:} Si $i^*$ no se conecta con ninguna piedra de su mismo color, evaluamos sus libertades. Si $L(i^*) = \varnothing$, la jugada constituye un suicidio.
		\item \textbf{Unión a un grupo:} Sea $I$ el grupo al que pertenece $i^*$ en este nuevo tablero. Si $L(I) = \varnothing$, la jugada  constituye un suicidio.
	\end{itemize}
	Si ocurre cualquiera de estos dos escenarios, la jugada es ilegal. Si no hay suicidio, avanzamos al siguiente paso.
	
	\item \textbf{Actualización:} Transformamos las piedras capturadas a estado libre ($Cap_0$) y actualizamos el tablero.
	$$Cap_0 = \{(0, P(i)) : i \in Cap\},\quad T_t = (T' \setminus Cap) \cup Cap_0.$$
\end{enumerate}
De esta manera, toda jugada propuesta puede aceptarse o rechazarse. Un Juego de Go se define formalmente como una cuádrupla,
$$ \Gamma = (G, E, A, \Delta).$$
Donde a un Go $(G, E, A)$, le adjuntamos la función de transición $\Delta: \mathcal{T} \times \mathcal{J} \to \mathcal{T}$.
%	Para ilustrar el método mediante la función emparejadora de Cantor, observemos primero el comportamiento del polinomio. Para el caso bidimensional ($n=2$), obtenemos:
%	\[ f(x_1,x_2)=\frac{(x_1 + x_2)^2 + 3x_1 + x_2}{2}.\]
%	Evaluando esta función para los 9 puntos del, obtenemos los siguientes valores:
%	\begin{table}[H]
	%		\centering
	%		\begin{tabular}{c c | c || c c | c}
		%			$x_1$ & $x_2$ & $f(x_1,x_2)$ & $x_1$ & $x_2$ & $f(x_1,x_2)$ \\
		%			\hline
		%			1 & 1 & 4  & 2 & 3 & 17 \\
		%			1 & 2 & 7  & 3 & 1 & 13 \\
		%			1 & 3 & 11 & 3 & 2 & 18 \\
		%			2 & 1 & 8  & 3 & 3 & 24 \\
		%			2 & 2 & 12 &   &   &    \\
		%		\end{tabular}
	%		\caption{Cómputo de la función de Cantor para un tablero de $3 \times 3$.}
	%		\label{tabla_cantor_3x3}
	%	\end{table}
%	Una vez calculados estos valores, construimos el vector de estados ordenando las posiciones de manera estrictamente creciente según su imagen bajo la función:
%	\[ T^{\pi} = \big( (e_1, 4), (e_2, 7), (e_3, 8), \dots, (e_9, 24) \big) \]
%	Este vector está en perfecta concordancia con el mapeo de las coordenadas originales ordenadas por su valor de Cantor:
%	\[ \tilde{T}^{\pi} = \big( (e_1, (1,1)), (e_2, (1,2)), (e_3, (2,1)), \dots, (e_9, (3,3)) \big) \]
%	Finalmente, si aplicamos nuestra función de estado $C$ sobre la secuencia ordenada $\tilde{T}^{\pi}$, aislamos exclusivamente la información de los colores, obteniendo el vector que usaremos para realizar las comparaciones y detectar el \textit{SuperKo}:
%	\[ C(\tilde{T}^{\pi}) = (e_1, e_2, \dots, e_9) = E^{\pi}. \]
%	
%	\vfill

%\section*{Glosario (Parte I)}
%\begin{center}
%	\begin{tabular}{p{2cm} p{4.7cm} p{2cm} p{4.7cm}}
	%		\toprule
	%		\textbf{Símbolo} & \textbf{Significado} & \textbf{Símbolo} & \textbf{Significado} \\
	%		\midrule
	%		$G$ & Puntos del tablero& $E$ & Estados $\{0,1,2\}$\\
	%		$C(i)$ & Estado de la piedra $i$ & $P(i)$ & Posición de la piedra $i$ \\
	%		$t_1\sim t_2$ & Adyacencia entre puntos & $i_1\!\leftrightarrow\!i_2$ & Conexión entre piedras \\
	%		$I^{e}$ & Grupo de piedras de color $e$ & $R$ & Grupo de puntos vacíos \\
	%		$K(I)$ & Adyacencias de $I$ & $L(I)$ & Libertades de $I$\\
	%		$J^{e}$ & Territorio del color $e$ & $f(t)$ & Jugada en el turno $t$\\
	%		 $H(t)$ & Hash de Zobrist del tablero en turno $t$ \\
	%		\bottomrule
	%	\end{tabular}
%\end{center}
%	\newpage
%	\twocolumn
%	\subsection{Resumen de la Parte I}
%	Un \textit{Go} es una $3$-upla, de la forma
%	$Go=(G,E,A)$.
%	$G=N_{1}\times N_{2},\dots, N_{n}$, $N_{i}\subset \mathbb{N}$, $n\in \mathbb{N}$. $|G|< \infty$
%	Un \textbf{Punto} es un elemento de $G$, notado por $t$. $A$ es la relación de adyacencia canónica, donde, dados dos \textbf{puntos} $t_{1}=(a_{1},\dots a_n)$, $t_{2}=(b_{1},\dots b_n)$ son \textbf{adyacentes}, notado por $t_{1}\sim t_{2}$, cuando $\exists! \ i \in \mathbb{N}$ tal que
%	\begin{align*}
	%	\forall r \in N_n \setminus\{i\},\ a_r=b_r  \ \land \  a_{i}=b_{i}\pm 1.
	%	\end{align*}
%	$E=\{0,1,\dots,m\}$, $0$ siempre reservado.\\
%	$E\times G=\{(e,t):e\in E, t\in G\}$.\\
%	Definimos $C$, $P$ por:\vspace{-5pt}
%	\begin{align*}
	%	C:E\times G \to E.&\,\, C(i)=C((e,t))=e.\\
	%	P:E\times G \to G.&\,\, P(i)=P((e,t))=t.
	%	\end{align*}
%	Una \textbf{piedra} es una $2$-upla, notada por $i$, de la forma $i=(e,t)$, para $e\in E, t \in G$.
%	Dos \textbf{piedras} $i_1=(e,t_{1})$ e $i_2=(e',t_{2})$, se dirá que son adyacentes cuando:\vspace{-5pt}
%	\begin{align*}
	%	i_1 \sim i_2 &\Leftrightarrow P(i_{1})\sim P(i_{2}). 
	%	\end{align*}
%	Decimos que $i_1$ e $i_2$ forman un \textbf{Grupo}, si sucede que,\vspace{-5pt}
%	\begin{align*}
	%		C(i_1)=C(i_2) &\land i_1 \sim i_2.
	%	\end{align*}
%	Se define la adyacencia de un grupo a un punto, denotado por $I_{r}(e)\sim t_{j}$. Cuando,
%	 $$\exists \ i_{k}=(e,t_{k}) \in I_{i}(e) : P(i_{k}) \sim t_{j}.$$
%	 $r \in \mathbb{N}$.
%	Definimos $R(i)$ y $R(I)$ por:
%	$$R(i_{j})=\{i_{r} \in E\times G: i_{j}\sim i_{r}\}.$$
%	$$R(I_{r}(e))=\{i_r: I_{i}(e)\sim i_{r}\}.$$
%	Una captura (cap) se da cuando, dado $i=(e,t), e\neq 0$, o $I_j(e)$: $0\notin C(R(i_{j}))$ o, $0\notin C(R(I_{j}(e)))$.
%	Luego se quitan las piedras, se establece, $C(i_{j})=0$ ó, $\forall i \in I_{j}(e), C(i)=0.$ Se dará un punto por cada piedra al jugador que realiza la captura, es decir $1$ punto si es una \textbf{piedra}, $|I_{j}(e)|$ puntos si se trata de un \textbf{grupo}.
%	Regla del suicidio, Sea $I_{j}(e)$, tal que, $0\in C(R(I_j(e)))$, si es turno de $e$ y, se juega de tal manera que $0\notin C(R(I_{j}(e)))$, dicha jugada es ilegal, pues se está auto-capturando es decir, suicidando. Si se juega donde hay \textbf{suicidio} y al mismo tiempo se \textbf{captura}, se considera una jugada legal.
%	Regla del SuperKo Posicional (SKo): Un jugador no puede colocar una piedra de manera tal que recree una posición que ya haya aparecido anteriormente en la partida.
%	Un \textbf{grupo} de color $0$, $I_{k}(0)$ se dice territorio del color $e\neq 0$, exactamente cuando:
%	\begin{align*}
	%	C(R(I_{j}(0))) &=\{e\}.
	%	\end{align*}
%	Los territorios de un color están dados por:
%	$$J^{e}=\{I_{r}(0): C(R(I_{r}(0)))=\{e\}\}.$$
%	Una jugada se realiza a partir de $f$, con:
%	$f:\mathbb{N}\rightarrow E\setminus\{0\}\times G \cup \varnothing$, tal que $f(n)= i_j=(e,t_j) \text{ ó } \varnothing$. No toda jugada es legal. $f(n) =i_{j}$ es legal si y solo si:
%	\begin{align*}
	%		C(i_{j})=0 \ \land \ \neg SKo(t_{j}) \ \land 
	%		\\ \ (\neg Suicidio(t_{j}) \lor Captura).
	%	\end{align*}
%	Entiéndase $\neg Sko$, como, no se ha roto la regla del \textbf{Sko}, no se ha repetido posición.
%	Para dar cuenta del \textbf{Sko}, hacemos uso del \textbf{Hash de Zobrist} \ref{ZobristHash}. Ya hemos discutido cómo dar cuenta del suicidio y las capturas.
%	
%	Si sucede que $f(n-m+1)=\dots=f(n-1)=f(n)=\varnothing$, entonces la partida termina. Quien gana en el \textbf{Go} es el jugador que posee más puntos, hacemos:
%	$$P(k)= \sum|I(k)|+	\sum|J(k)|+ Cap(k)$$
%	Para cada $k \in E\setminus \{0\}$, el mayor $P(k)$ gana. Si hay hay más de uno que es mayor, es empate. Aquí no consideramos \textbf{komi} por ser el generalizado, en cuyo caso (dos jugadores), sumar \textbf{komi} al jugador $2$.
%	
%	Un $Go$ además es notado por $Go^{n}_{m}$. Si quisiéramos explicitar: $Go^{2}_{2}(19,19)$ para el tablero convencional. En caso de omitir los jugadores, se interpretará que se está hablando del tablero de dimensión $n$, para cualquier cantidad de jugadores: $Go^{n}$. Del mismo modo, $Go_{m}$ representará el tablero de cualquier dimensión y de $m$ jugadores.
%\onecolumn
%	\newpage
\section{Vida en el juego del Go}
\label{sec-vida}
Si recordamos las reglas, justo hacia el final habíamos visto brevemente el concepto de \textit{vivo}. En esta sección lo formalizaremos.


\begin{wrapfigure}[12]{r}{0.3\textwidth}
	%\begin{figure}[H]
	\centering
	\begin{tikzpicture}[scale=0.485] \tikzstyle{every node}=[font=\small]
		% Draw the background
		\fill[goban] (-0.5,-0.5) rectangle (6.5,6.5);
		
		% Draw the grid lines
		\draw[step=1cm, black] (0,0) grid (6,6);
		
		% Draw the thick border
		\draw[very thick] (0,0) rectangle (6,6);
		\filldraw[draw=black,fill=white] (2,1) circle (12pt);
		\filldraw[draw=black,fill=white] (2,2) circle (12pt);
		\filldraw[draw=black,fill=white] (0,2) circle (12pt);
		\filldraw[draw=black,fill=white] (1,2) circle (12pt);
		\filldraw[draw=black,fill=white] (3,2) circle (12pt);
		\filldraw[draw=black,fill=white] (4,2) circle (12pt);
		\filldraw[draw=black,fill=white] (4,1) circle (12pt);
		\filldraw[draw=black,fill=white] (4,0) circle (12pt);
		\filldraw[black] (0,1) circle (12pt);
		\filldraw[black] (1,1) circle (12pt);
		\filldraw[black] (1,0) circle (12pt);
		\filldraw[black] (3,1) circle (12pt);
		\filldraw[black] (2,1) circle (12pt);
		\filldraw[black] (3,0) circle (12pt);
		% Draw the central black dot
	\end{tikzpicture}
	\caption{Tablero con un grupo negro que posee vida incondicional.}
	\label{fig:vida_incond_ejemplo}
\end{wrapfigure}En la Figura~\ref{fig:vida_incond_ejemplo}, el blanco no puede capturar porque la regla del suicidio le impide jugar en las libertades restantes del negro. El algoritmo para determinar la vida incondicional es sencillo: se simula que el negro pasa mientras el blanco intenta ocupar sus libertades. Si el blanco no logra capturarlo, el grupo tiene vida incondicional.	En 1976, David B. Benson formalizó este concepto \cite{Benson}. Demostró que la supervivencia de un conjunto de piedras depende únicamente de sus adyacencias. A continuación, adaptaremos sus definiciones y teoremas a nuestro marco formal.

Para llegar a ello, primero necesitamos contar cuántos grupos rodean a una región. Sea $P$ un conjunto de piedras de color $e \in E \setminus \{0\}$ y $R$ una región rodeada completamente por $P$ (es decir, $K(R) \subseteq P$). Nos interesa saber, conociendo únicamente $R$, cuántos grupos distintos componen su \textit{frontera} en $P$.	Sabemos que el conjunto $K(R)$ contiene las piedras adyacentes a la región. Por lo visto en la Parte~\ref{parte1}, $0 \notin C(K(R))$.
\begin{definition}[Frontera de una región]
	Sea $R$ una región. Llamaremos frontera de $R$, denotada como $\partial R$, al conjunto de todos los grupos o piedras que tienen una intersección no nula con las adyacencias de la región.
	$$ \partial R = \{ I \mid I \cap K(R) \neq \varnothing \}\cup \{ i \mid i \cap K(R) \neq \varnothing \}.$$
\end{definition}
Es crucial notar la diferencia: $\partial R$ agrupa los grupos adyacentes, mientras que $K(R)$ agrupa las piedras adyacentes, la Figura~\ref{fig:dif_k_partial} ilustra esta distinción.
\begin{figure}[H]
	\centering
	\begin{tikzpicture}[scale=0.485] \tikzstyle{every node}=[font=\small]
		% Draw the background
		\fill[goban] (-0.5,-0.5) rectangle (6.5,6.5);
		
		% Draw the grid lines
		\draw[step=1cm, black] (0,0) grid (6,6);
		
		% Draw the thick border
		\draw[very thick] (0,0) rectangle (6,6);
		
		\filldraw[black] (3,2) circle (12pt);
		\filldraw[black] (4,2) circle (12pt);
		\filldraw[OliveGreen] (3,2) circle (4pt);
		\filldraw[OliveGreen] (4,2) circle (4pt);
		
		\filldraw[draw=black,fill=white] (2,3) circle (12pt);
		\filldraw[draw=black,fill=white] (1,3) circle (12pt);
		\filldraw[draw=black,fill=white] (0,3) circle (12pt);
		\filldraw[magenta] (2,3) circle (4pt);
		\filldraw[magenta] (1,3) circle (4pt);
		\filldraw[magenta] (0,3) circle (4pt);
		
		\filldraw[draw=black,fill=white] (3,4) circle (12pt);
		\filldraw[draw=black,fill=white] (4,4) circle (12pt);
		\filldraw[draw=black,fill=white] (3,5) circle (12pt);
		\filldraw[draw=black,fill=white] (2,5) circle (12pt);
		\filldraw[draw=black,fill=white] (4,4) circle (12pt);
		\filldraw[draw=black,fill=white] (5,3) circle (12pt);
		\filldraw[draw=black,fill=white] (6,2) circle (12pt);
		\filldraw[Orange] (3,4) circle (4pt);
		\filldraw[Orange] (3,5) circle (4pt);
		\filldraw[Orange] (2,5) circle (4pt);
		\filldraw[Orange] (4,4) circle (4pt);
		\filldraw[Blue] (5,3) circle (4pt);
		\fill[draw=black,fill=Red, very thin] (3.25,3.25) rectangle (2.75,2.75);
		\fill[draw=black,fill=Red, very thin] (4.25,3.25) rectangle (3.75,2.75);
		
		\begin{scope}[xshift=8cm]
			\fill[goban] (-0.5,-0.5) rectangle (6.5,6.5);
			
			% Draw the grid lines
			\draw[step=1cm, black] (0,0) grid (6,6);
			
			% Draw the thick border
			\draw[very thick] (0,0) rectangle (6,6);
			
			\filldraw[black] (3,2) circle (12pt);
			\filldraw[black] (4,2) circle (12pt);
			\filldraw[OliveGreen] (3,2) circle (4pt);
			\filldraw[OliveGreen] (4,2) circle (4pt);
			\filldraw[draw=black,fill=white] (2,3) circle (12pt);
			\filldraw[draw=black,fill=white] (1,3) circle (12pt);
			\filldraw[draw=black,fill=white] (0,3) circle (12pt);
			\filldraw[OliveGreen] (2,3) circle (4pt);
			
			\filldraw[draw=black,fill=white] (3,4) circle (12pt);
			\filldraw[draw=black,fill=white] (4,4) circle (12pt);
			\filldraw[draw=black,fill=white] (3,5) circle (12pt);
			\filldraw[draw=black,fill=white] (2,5) circle (12pt);
			\filldraw[draw=black,fill=white] (4,4) circle (12pt);
			\filldraw[draw=black,fill=white] (5,3) circle (12pt);
			\filldraw[draw=black,fill=white] (6,2) circle (12pt);
			\filldraw[OliveGreen] (3,4) circle (4pt);
			\filldraw[OliveGreen] (4,4) circle (4pt);
			\filldraw[OliveGreen] (5,3) circle (4pt);
			\fill[draw=black,fill=Red, very thin] (3.25,3.25) rectangle (2.75,2.75);
			\fill[draw=black,fill=Red, very thin] (4.25,3.25) rectangle (3.75,2.75);
		\end{scope}
		
		% Draw the central black dot
	\end{tikzpicture}
	\caption{Izquierda. Tablero con los \textit{grupos frontera} de una \textcolor{Red}{Región}; \textcolor{Orange}{Naranja}, \textcolor{magenta}{Magenta}, \textcolor{OliveGreen}{Verde}, \textcolor{Blue}{Azul}. Derecha. \textcolor{OliveGreen}{Adyacencias} de una \textcolor{Red}{Región}.}
	\label{fig:dif_k_partial}
\end{figure}

El cardinal de este conjunto nos indica exactamente cuántos grupos y piedras independientes rodean a la región vacía. 
%	\begin{figure}[H]
	%		\centering
	%		\begin{tikzpicture}[scale=0.485] \tikzstyle{every node}=[font=\small]
		%			% Draw the background
		%			\fill[goban] (-0.5,-0.5) rectangle (6.5,6.5);
		%			
		%			% Draw the grid lines
		%			\draw[step=1cm, black] (0,0) grid (6,6);
		%			
		%			% Draw the thick border
		%			\draw[very thick] (0,0) rectangle (6,6);
		%			\filldraw[black] (0,1) circle (12pt);
		%			\filldraw[black] (1,0) circle (12pt);
		%			\filldraw[Salmon] (0,1) circle (4pt);
		%			\filldraw[Orange] (1,0) circle (4pt);
		%			
		%			\filldraw[black] (2,2) circle (12pt);
		%			\filldraw[magenta] (2,2) circle (4pt);
		%			\filldraw[black] (3,3) circle (12pt);
		%			\filldraw[magenta] (3,3) circle (4pt);
		%			\filldraw[black] (4,3) circle (12pt);
		%			\filldraw[magenta] (4,3) circle (4pt);
		%			\filldraw[black] (3,1) circle (12pt);
		%			\filldraw[Violet] (3,1) circle (4pt);
		%			\filldraw[black] (2,3) circle (12pt);
		%			\filldraw[Magenta] (2,3) circle (4pt);
		%			\filldraw[draw=black,fill=white] (4,1) circle (12pt);
		%			\filldraw[OliveGreen] (4,1) circle (4pt);
		%			\filldraw[draw=black,fill=white] (5,2) circle (12pt);
		%			\filldraw[RawSienna] (5,2) circle (4pt);
		%
		%			\fill[draw=black,fill=Blue, very thin] (3.25,1.75) rectangle (2.75,2.25);
		%			\fill[draw=black,fill=Blue, very thin] (4.25,1.75) rectangle (3.75,2.25);
		%			\fill[draw=black,fill=Red, very thin] (0.25,0.25) rectangle (-0.25,-0.25);
		%
		%			% Draw the central black dot
		%		\end{tikzpicture}
	%		\caption{Tablero con un grupos frontera; \textcolor{Orange}{Naranja}, \textcolor{Salmon}{Salmón}; \textcolor{magenta}{Magenta}, \textcolor{Violet}{Violeta}, \textcolor{OliveGreen}{Verde}, \textcolor{RawSienna}{Marrón}. Que rodean respectivamente a una región \textcolor{Red}{Roja} y otra \textcolor{Blue}{Azul}.}
	%		\label{fig:ej_frontera}
	%	\end{figure}
Sin embargo, no basta con la cantidad; también debemos evaluar la calidad de estas regiones. En términos de Benson, deben ser \textit{saludables} para los grupos que las contienen.

\begin{wrapfigure}[10]{r}{0.3\textwidth}
	\vspace{-10pt}
	\centering
	\begin{tikzpicture}[scale=0.485] \tikzstyle{every node}=[font=\small]
		% Draw the background
		\fill[goban] (-0.5,-0.5) rectangle (6.5,6.5);
		
		% Draw the grid lines
		\draw[step=1cm, black] (0,0) grid (6,6);
		
		% Draw the thick border
		\draw[very thick] (0,0) rectangle (6,6);
		\filldraw[draw=black,fill=white] (0,0) circle (12pt);
		\filldraw[draw=black,fill=white] (0,1) circle (12pt);
		\filldraw[draw=black,fill=white] (0,2) circle (12pt);
		\filldraw[draw=black,fill=white] (1,2) circle (12pt);
		\filldraw[draw=black,fill=white] (2,3) circle (12pt);
		\filldraw[draw=black,fill=white] (3,3) circle (12pt);
		\filldraw[draw=black,fill=white] (4,3) circle (12pt);
		\filldraw[draw=black,fill=white] (4,2) circle (12pt);
		\filldraw[draw=black,fill=white] (5,2) circle (12pt);
		\filldraw[draw=black,fill=white] (6,2) circle (12pt);
		\filldraw[draw=black,fill=white] (6,1) circle (12pt);
		\filldraw[draw=black,fill=white] (6,0) circle (12pt);
		\filldraw[black] (1,0) circle (12pt);
		\filldraw[black] (1,1) circle (12pt);
		\filldraw[black] (2,2) circle (12pt);
		\filldraw[black] (3,2) circle (12pt);
		\filldraw[black] (3,1) circle (12pt);
		\filldraw[black] (4,1) circle (12pt);
		\filldraw[black] (5,1) circle (12pt);
		\filldraw[black] (5,0) circle (12pt);
		
		\fill[draw=black,fill=OliveGreen, very thin] (2.25,0.25) rectangle (1.75,-0.25);
		\fill[draw=black,fill=OliveGreen, very thin] (3.25,0.25) rectangle (2.75,-0.25);
	\end{tikzpicture}
	\caption{Tablero con grupos negros cuya \textcolor{Blue}{región} controlada es \textcolor{OliveGreen}{no saludable}.}
	\label{fig:benson_ej_2}
\end{wrapfigure}\begin{definition}[Región saludable]
	Sea $P$ un conjunto de grupos y piedras de un mismo color y $R$ una región rodeada por $P$ (es decir, $\partial R \subset P$). Se dice que $R$ es saludable para un grupo $I \in P$ (o $i\in P$) si todo elemento de $R$ es una libertad para $I$ ($i$). Formalmente:
	$$R \text{ es saludable para } I \in P \iff R \subseteq L(I).$$
\end{definition}Esta condición asegura que, si el oponente intenta jugar dentro de la región, quedará inmediatamente en contacto con el grupo $I$ ($i$). En el ejemplo de la Figura~\ref{fig:benson_ej_2}, Los puntos marcados en \textcolor{OliveGreen}{verde} no cumplen con esta premisa; por lo tanto, esa región no es saludable para ningún grupo negro. \break

\begin{definition}[Región vital]\label{def:reg_vital}
	Sea $P$ un conjunto de grupos y piedras de un mismo color, $I \in P$ un grupo específico y $R$ una región. Diremos que $R$ es vital para $I$ en $P$, si se cumplen simultáneamente las siguientes condiciones:
	\begin{enumerate}
		\item $R$ es saludable para $I$, $R\subseteq L(I)$.
		\item Todos los grupos que conforman la frontera de $R$ pertenecen a $P$, $\partial R \subseteq P$.
	\end{enumerate}
\end{definition}
A las regiones vitales las notaremos por $V$. Esta definición formaliza que una región debe proveer libertades inmediatas (ser saludable) y estar rodeada por grupos de un mismo color. A partir de la noción de región vital, podemos clasificar la naturaleza de las regiones completamente rodeadas.
\begin{definition}[Ojo]
	Sea $V$ una región vital para un grupo $I$ (o piedra) dentro de un conjunto $P$. Si la frontera de $V$ está conformada por un único grupo independiente, es decir, si $|\partial V| = 1$, entonces $V$ será un ojo para $I$.
\end{definition}
Cuando la región está rodeada por múltiples grupos distintos, del mismo color, estamos ante una \textit{unión}.
\begin{definition}[Unión]
	Sea $R$ una región saludable para un grupo $I$ (o piedra) dentro de un conjunto $P$. Si la frontera de $R$ está compuesta por más de un grupo independiente, es decir, si $|\partial R| > 1$. Entonces $R$ es una unión para el conjunto $P$ («\textit{joiner}» para Benson). La Figura~\ref{fig:benson_ejemplos} ilustra la sutil diferencia entre las definiciones de Benson.
\end{definition}
\begin{figure}[H]
	\centering
	\begin{tikzpicture}[scale=0.485] \tikzstyle{every node}=[font=\small]
		% --------------------------------------------------------
		% TABLERO 1: Saludable vs No Saludable
		% --------------------------------------------------------
		\fill[goban] (-0.5,-0.5) rectangle (6.5,6.5);
		\draw[step=1cm, black] (0,0) grid (6,6);
		\draw[very thick] (0,0) rectangle (6,6);
		
		% Draw the thick border
		\draw[very thick] (0,0) rectangle (6,6);
		\filldraw[draw=black,fill=white] (0,0) circle (12pt);
		\filldraw[draw=black,fill=white] (0,1) circle (12pt);
		\filldraw[draw=black,fill=white] (0,2) circle (12pt);
		\filldraw[draw=black,fill=white] (1,2) circle (12pt);
		\filldraw[draw=black,fill=white] (2,3) circle (12pt);
		\filldraw[draw=black,fill=white] (3,3) circle (12pt);
		\filldraw[draw=black,fill=white] (4,3) circle (12pt);
		\filldraw[draw=black,fill=white] (4,2) circle (12pt);
		\filldraw[draw=black,fill=white] (5,2) circle (12pt);
		\filldraw[draw=black,fill=white] (6,2) circle (12pt);
		\filldraw[draw=black,fill=white] (6,1) circle (12pt);
		\filldraw[draw=black,fill=white] (6,0) circle (12pt);
		\filldraw[draw=black,fill=white] (1,3) circle (12pt);
		\filldraw[black] (1,0) circle (12pt);
		\filldraw[black] (1,1) circle (12pt);
		\filldraw[black] (2,2) circle (12pt);
		\filldraw[black] (3,2) circle (12pt);
		\filldraw[black] (3,1) circle (12pt);
		\filldraw[black] (4,1) circle (12pt);
		\filldraw[black] (5,1) circle (12pt);
		\filldraw[black] (5,0) circle (12pt);
		
		\fill[draw=black,fill=OliveGreen, very thin] (2.25,0.25) rectangle (1.75,-0.25);
		\fill[draw=black,fill=OliveGreen, very thin] (3.25,0.25) rectangle (2.75,-0.25);
		%				\fill[draw=black,fill=OliveGreen, very thin] (4.25,0.25) rectangle (3.75,-0.25);
		%				\fill[draw=black,fill=OliveGreen, very thin] (2.25,1.25) rectangle (1.75,0.75);
		
		\filldraw[black] (0,4) circle (12pt);
		\filldraw[black] (1,4) circle (12pt);
		\filldraw[black] (2,4) circle (12pt);
		\filldraw[black] (3,4) circle (12pt);
		\filldraw[black] (3,5) circle (12pt);
		\filldraw[black] (0,6) circle (12pt);
		\filldraw[black] (1,6) circle (12pt);
		\filldraw[black] (2,6) circle (12pt);
		
		\fill[draw=black,fill=Orange, very thin] (-0.25,4.75) rectangle (0.25,5.25);
		\fill[draw=black,fill=Orange, very thin] (0.75,4.75) rectangle (1.25,5.25);
		\fill[draw=black,fill=Orange, very thin] (1.75,4.75) rectangle (2.25,5.25);
		% --------------------------------------------------------
		% TABLERO 2: ojo vs union
		% --------------------------------------------------------
		\begin{scope}[xshift=8cm]
			\fill[goban] (-0.5,-0.5) rectangle (6.5,6.5);
			\draw[step=1cm, black] (0,0) grid (6,6);
			\draw[very thick] (0,0) rectangle (6,6);
			
			\filldraw[black] (3,6) circle (12pt);
			\filldraw[black] (0,4) circle (12pt);
			\filldraw[black] (1,4) circle (12pt);
			\filldraw[black] (2,4) circle (12pt);
			\filldraw[black] (3,4) circle (12pt);
			\filldraw[black] (3,5) circle (12pt);
			\filldraw[black] (0,6) circle (12pt);
			\filldraw[black] (1,6) circle (12pt);
			\filldraw[black] (2,6) circle (12pt);
			\fill[draw=black,fill=purple, very thin] (-0.25,4.75) rectangle (0.25,5.25);
			\fill[draw=black,fill=purple, very thin] (0.75,4.75) rectangle (1.25,5.25);
			\fill[draw=black,fill=purple, very thin] (1.75,4.75) rectangle (2.25,5.25);
			
			\filldraw[draw=black,fill=white] (2,1) circle (12pt);
			\filldraw[draw=black,fill=white] (2,2) circle (12pt);
			\filldraw[draw=black,fill=white] (0,2) circle (12pt);
			\filldraw[draw=black,fill=white] (1,2) circle (12pt);
			\filldraw[draw=black,fill=white] (3,2) circle (12pt);
			\filldraw[draw=black,fill=white] (4,2) circle (12pt);
			\filldraw[draw=black,fill=white] (4,1) circle (12pt);
			\filldraw[draw=black,fill=white] (3,1) circle (12pt);
			\filldraw[black] (0,1) circle (12pt);
			\filldraw[black] (1,1) circle (12pt);
			\filldraw[black] (1,0) circle (12pt);
			\filldraw[black] (3,1) circle (12pt);
			\filldraw[draw=black,fill=white] (2,1) circle (12pt);
			\filldraw[draw=black,fill=white] (5,3) circle (12pt);
			\filldraw[black] (4,1) circle (12pt);
			\filldraw[black] (4,0) circle (12pt);
			\filldraw[black] (2,0) circle (12pt);
			\filldraw[black] (5,2) circle (12pt);
			\filldraw[black] (6,2) circle (12pt);
			\filldraw[black] (6,1) circle (12pt);
			\filldraw[black] (6,0) circle (12pt);
			\fill[draw=black,fill=purple, very thin] (-0.25,-0.25) rectangle (0.25,0.25);
			\fill[draw=black,fill=cyan, very thin] (2.75,-0.25) rectangle (3.25,0.25);
			\fill[draw=black,fill=cyan, very thin] (4.75,0.75) rectangle (5.25,1.25);
			\fill[draw=black,fill=cyan, very thin] (4.75,-0.25) rectangle (5.25,0.25);
			
		\end{scope}
		
	\end{tikzpicture}
	\caption{Izquierda: \textcolor{Orange}{saludable} y \textcolor{OliveGreen}{no saludable}. Derecha: \textcolor{purple}{ojos} y \textcolor{cyan}{uniones}.}
	\label{fig:benson_ejemplos}
\end{figure}

\subsection{Teorema de Benson}
Apliquemos la teoría y veamos el Teorema de Benson. Antes de llegar a él, necesitaremos establecer cuatro lemas previos. En lo que sigue, supondremos que el jugador que defiende el grupo atacado pasa en cada turno.
\begin{lemma}\label{lem_1_benson}
	Si un grupo $I$ posee al menos dos ojos distintos, entonces tiene vida incondicional.
\end{lemma}
\begin{proof}
	Sean $O_1$ y $O_2$ los ojos de $I$. Para capturar al grupo $I$, el oponente debe ocupar todas sus libertades, lo que implica llenar completamente $O_1$ y $O_2$.
	
	Sin pérdida de generalidad, supongamos que el oponente intenta llenar primero $O_1$. Dado que la región es finita, eventualmente quedará un único punto $x \in O_1$.
	
	Al jugar en $x$, las piedras del oponente dentro de $O_1$ se quedan sin libertades, ya que $O_1$ está rodeado exclusivamente por $I$ ($|\partial O_1| = 1$). Dado que $I$ aún posee al menos una libertad en $O_2$, la jugada en $x$ no captura a $I$. Por consiguiente, la regla del suicidio determina que ocupar el punto $x$ es ilegal.
	
	Como resultado, es imposible eliminar por completo ninguno de los dos ojos mediante jugadas legales. Por lo tanto, $I$ tiene vida incondicional.
\end{proof}
%	\begin{figure}[H]
	%		\centering
	%		\begin{tikzpicture}[scale=0.485] \tikzstyle{every node}=[font=\small]
		%			% Draw the background
		%			\fill[goban] (-0.5,-0.5) rectangle (6.5,6.5);
		%			
		%			% Draw the grid lines
		%			\draw[step=1cm, black] (0,0) grid (6,6);
		%			
		%			% Draw the thick border
		%			\draw[very thick] (0,0) rectangle (6,6);
		%			\filldraw[draw=black,fill=white] (2,1) circle (12pt);
		%			\filldraw[draw=black,fill=white] (2,2) circle (12pt);
		%	
		%
		%			\filldraw[draw=black,fill=white] (3,2) circle (12pt);
		%			\filldraw[draw=black,fill=white] (4,2) circle (12pt);
		%			\filldraw[draw=black,fill=white] (4,1) circle (12pt);
		%			\filldraw[draw=black,fill=white] (3,1) circle (12pt);
		%			\filldraw[black] (0,1) circle (12pt);
		%			\filldraw[black] (1,1) circle (12pt);
		%			\filldraw[black] (1,0) circle (12pt);
		%			\filldraw[black] (3,1) circle (12pt);
		%			\filldraw[black] (2,1) circle (12pt);
		%			\filldraw[black] (5,0) circle (12pt);
		%			\filldraw[black] (5,1) circle (12pt);
		%			\filldraw[black] (1,2) circle (12pt);
		%			\filldraw[black] (0,3) circle (12pt);
		%			\filldraw[draw=black,fill=white] (0,4) circle (12pt);
		%			\filldraw[draw=black,fill=white] (1,3) circle (12pt);
		%			\fill[draw=black,fill=Red, very thin] (0.25,0.25) rectangle (-0.25,-0.25);
		%			\fill[draw=black,fill=OliveGreen, very thin] (0.25,2.25) rectangle (-0.25,1.75);
		%			% Draw the central black dot
		%		\end{tikzpicture}
	%		\caption{Tablero con un \textcolor{Red}{Ojo} y una \textcolor{OliveGreen}{Unión}.}
	%		\label{fig:dif_Unión}
	%	\end{figure}
\begin{lemma}\label{lem_2_benson}
	Sea $S$ un conjunto de grupos de un mismo color tal que $|S| \ge 2$. Si un grupo $I_1 \in S$ posee un ojo $O$ y al menos una unión $U$ en cuya frontera se encuentra otro grupo distinto $I_2 \in S$, entonces $I_2$ debe ser capturado antes de poder capturar a $I_1$.
\end{lemma}
\begin{proof}
	Para capturar a $I_{1}$, el oponente debe ocupar todas las libertades provistas por la unión $U$, pues las del ojo $O$ solo pueden llenarse en última instancia (Lema~\ref{lem_1_benson}).
	
	Dado que $U$ es una unión, su frontera contiene más de un grupo ($|\partial U| > 1$), incluyendo a $I_1$ e $I_2$. Toda jugada en $U$ disminuye simultáneamente las libertades de ambos. Al intentar ocupar la última libertad vacía de $U$, el movimiento del oponente carecerá de libertades propias. Por la regla del suicidio, esta jugada es legal solo si realiza una captura.
	
	Como $I_1$ conserva al menos una libertad en su ojo $O$, no puede ser capturado en este turno. Por lo tanto, la jugada debe capturar obligatoriamente a otro grupo en la frontera de $U$, en este caso $I_2$. En consecuencia, $I_2$ es capturado antes que $I_1$.
\end{proof}
En síntesis, para capturar a un grupo que posee un ojo y está conectado a una unión, es necesario capturar primero a los demás grupos que forman parte de la frontera de dicha unión, como se ilustra en la Figura~\ref{fig:lem_benson}.
\begin{lemma}\label{lem_3_benson}
	Sea $I$ un grupo sin ojos, pero conectado a $n$ uniones $U_1, U_2, \dots, U_n$ ($n \ge 2$). Entonces, antes de poder capturar a $I$, el oponente debe capturar al menos un grupo distinto en $n-1$ de estas uniones.
\end{lemma}
\begin{proof}
	Para capturar a $I$, el oponente debe eliminar todas sus libertades en las $n$ uniones. Como no puede ocupar simultáneamente la última libertad de todas ellas, la regla del suicidio le obliga a reservar al menos una unión, digamos $U_n$, para la jugada final.
	
	Para agotar las libertades en las $n-1$ uniones restantes ($U_1, \dots, U_{n-1}$), cada unión presenta la situación descrita en el Lema~\ref{lem_2_benson}. Por lo tanto, el oponente se verá forzado a capturar al menos un grupo distinto en cada una de esas $n-1$ uniones antes de poder atacar a $I$.
\end{proof}
\begin{figure}[H]
	\centering
	\begin{tikzpicture}[scale=0.485] \tikzstyle{every node}=[font=\small]
		% Draw the background
		\fill[goban] (-0.5,-0.5) rectangle (6.5,6.5);
		
		% Draw the grid lines
		\draw[step=1cm, black] (0,0) grid (6,6);
		
		% Draw the thick border
		\draw[very thick] (0,0) rectangle (6,6);
		\filldraw[draw=black,fill=white] (2,1) circle (12pt);
		\filldraw[draw=black,fill=white] (2,2) circle (12pt);
		\filldraw[draw=black,fill=white] (0,2) circle (12pt);
		\filldraw[draw=black,fill=white] (1,2) circle (12pt);
		\filldraw[draw=black,fill=white] (3,2) circle (12pt);
		\filldraw[draw=black,fill=white] (4,2) circle (12pt);
		\filldraw[draw=black,fill=white] (4,1) circle (12pt);
		\filldraw[draw=black,fill=white] (3,1) circle (12pt);
		\filldraw[black] (0,1) circle (12pt);
		\filldraw[black] (1,1) circle (12pt);
		\filldraw[black] (1,0) circle (12pt);
		\filldraw[black] (3,1) circle (12pt);
		\filldraw[draw=black,fill=white] (2,1) circle (12pt);
		\filldraw[draw=black,fill=white] (5,1) circle (12pt);
		\filldraw[black] (4,1) circle (12pt);
		\filldraw[black] (4,0) circle (12pt);
		\filldraw[black] (2,0) circle (12pt);
		\filldraw[Red] (0,0) circle (6pt);
		\filldraw[Blue] (3,0) circle (6pt);
		\fill[draw=black,fill=Red, very thin] (-0.25,-0.25) rectangle (0.25,0.25);
		\fill[draw=black,fill=Blue, very thin] (2.75,-0.25) rectangle (3.25,0.25);
		
		\begin{scope}[xshift=8cm]
			\fill[goban] (-0.5,-0.5) rectangle (6.5,6.5);
			
			% Draw the grid lines
			\draw[step=1cm, black] (0,0) grid (6,6);
			
			% Draw the thick border
			\draw[very thick] (0,0) rectangle (6,6);
			\filldraw[draw=black,fill=white] (1,3) circle (12pt);
			\filldraw[draw=black,fill=white] (1,6) circle (12pt);
			\filldraw[draw=black,fill=white] (1,5) circle (12pt);
			\filldraw[draw=black,fill=white] (2,3) circle (12pt);
			\filldraw[draw=black,fill=white] (2,5) circle (12pt);
			\filldraw[draw=black,fill=white] (1,4) circle (12pt);
			\filldraw[draw=black,fill=white] (3,4) circle (12pt);
			\filldraw[draw=black,fill=white] (3,6) circle (12pt);
			\filldraw[draw=black,fill=white] (4,6) circle (12pt);
			\filldraw[draw=black,fill=white] (5,6) circle (12pt);
			\filldraw[draw=black,fill=white] (6,5) circle (12pt);
			\filldraw[black] (0,6) circle (12pt);
			\filldraw[black] (0,5) circle (12pt);
			\filldraw[black] (0,4) circle (12pt);
			\filldraw[black] (0,3) circle (12pt);
			\filldraw[black] (0,2) circle (12pt);
			\filldraw[black] (1,2) circle (12pt);
			\filldraw[black] (2,2) circle (12pt);
			\filldraw[black] (3,2) circle (12pt);
			\filldraw[black] (4,2) circle (12pt);
			\filldraw[black] (4,3) circle (12pt);
			\filldraw[black] (4,4) circle (12pt);
			\filldraw[black] (4,5) circle (12pt);
			\filldraw[black] (5,5) circle (12pt);
			\fill[draw=black,fill=Blue, very thin] (1.75,5.75) rectangle (2.25,6.25);
			\fill[draw=black,fill=Blue, very thin] (5.75,5.75) rectangle (6.25,6.25);
			\fill[draw=black,fill=Blue, very thin] (1.75,3.75) rectangle (2.25,4.25);
		\end{scope}
		% Draw the central black dot
	\end{tikzpicture}
	\caption{Izquierda: El grupo negro posee un ojo (rojo) y una unión (azul); el blanco debe capturar primero al grupo derecho (Lema~\ref{lem_2_benson}). Derecha: El grupo blanco superior izquierdo carece de ojos, pero comparte tres uniones (azul); el negro debe capturar a los grupos vecinos en al menos dos de ellas antes de poder capturarlo (Lema~\ref{lem_3_benson}).}
	\label{fig:lem_benson}
\end{figure}
\begin{lemma}\label{lem_4_benson}
	Sea $I$ un grupo o piedra, las libertades que no pertenecen a regiones vitales pueden ocuparse una a una mediante una secuencia de jugadas legales.
\end{lemma}
\begin{proof}
	Sea $L(I)$ el conjunto de libertades de $I$. Toda libertad $x \in L(I)$ pertenece a alguna región vacía $R$. Si $R$ no es una región vital para el conjunto $P$ (al cual pertenece $I$), significa que incumple la segunda condición de la Definición~\ref{def:reg_vital}; es decir, su frontera no está contenida enteramente en $P$ ($\partial R \not\subset P$).
	
	Al existir elementos en la frontera de $R$ que no pertenecen a $P$, el oponente puede colocar piedras en $R$ sin que estas se queden sin libertades. Como $\partial R \not\subset P$, las piedras enemigas jugadas en la región mantendrán o encontrarán libertades a través de esos grupos externos a $P$, evitando que la jugada sea un suicidio. Por lo tanto, el oponente puede ocupar estas libertades una a una de manera legal.
\end{proof}
\begin{theorem}[Vida incondicional de Benson]\label{teo_benson}
	Sea $I$ un grupo. Bajo un juego de $\mathrm{Go}_{m}$ canónico, ($m \geq 2$). Entonces, $I$ posee vida incondicional si y solo si pertenece a un conjunto de grupos $S$ de un mismo color tal que, para cada grupo $J \in S$, existen al menos dos regiones vitales distintas en $S$.
\end{theorem}
Comenzaremos demostrando la vuelta. Luego, mediante el contrarrecíproco, la ida.
\begin{proof}
	$(\Leftarrow):$ Sea $I\in S$. Para capturarlo, debemos rodearlo completamente. En particular, debemos ocupar las libertades de sus regiones vitales. Por los Lemas~\ref{lem_1_benson}, \ref{lem_2_benson}, y \ref{lem_3_benson}, ya sean estas regiones ojos o uniones, la captura de $I$ exige la captura previa de algún otro grupo de $S$ que forme parte de la frontera de dichas regiones.
	
	Llamemos $I_k \in S$ a ese grupo. Para capturarlo, eventualmente debemos ocupar las libertades de sus propias regiones vitales (que podrían coincidir con las de $I$). Para ello, nuevamente, debemos capturar a otro grupo que pertenezca a la frontera de sus regiones vitales.
	
	Repitiendo el argumento, obtenemos que toda captura de un grupo de $S$ requiere una captura estrictamente previa de otro grupo de $S$. Por consiguiente, ninguna secuencia finita de jugadas puede contener una primera captura de un grupo de $S$. Por lo tanto, $S$ posee vida incondicional y, consecuentemente, $I$ también.
	
	$(\Rightarrow):$ Por contrarrecíproco, demostraremos que:
	$$ \exists I \in S \text{ con a lo sumo una región vital en } S \implies S \text{ no posee vida incondicional.} $$
	
	Sea $I \in S$ un grupo que posee, a lo sumo, una única región vital $V$ en $S$, por el lema~\ref{lem_4_benson}, mediante una secuencia de jugadas legales es posible ocupar todas las libertades externas de $I$. Una vez ocupadas las libertades externas, analizamos el estado de $I$:
	\begin{itemize}
		\item Si $I$ carecía de regiones vitales, la ocupación del exterior agota sus libertades, concretando su captura.
		\item Si $I$ poseía exactamente una región vital $V$, todas sus libertades están en $V$. Al intentar ocupar el último punto de $V$, como $I$ no posee una segunda región vital que le garantice al menos una libertad externa a la jugada, colocar esta piedra elimina la ultima libertad de $I$, efectuando la captura.
	\end{itemize}
	
	En ambos casos, existe una secuencia válida para reducir las libertades de $I$ y capturarlo. Por lo tanto, $S$ no posee vida incondicional, concluyendo la demostración.
\end{proof}

\begin{wrapfigure}[8]{r}{0.3\textwidth}
	\vspace{-10pt}
	\centering
	\begin{tikzpicture}[scale=0.485] \tikzstyle{every node}=[font=\small]
		% Draw the background
		\fill[goban] (-0.5,-0.5) rectangle (6.5,6.5);
		
		% Draw the grid lines
		\draw[step=1cm, black] (0,0) grid (6,6);
		
		% Draw the thick border
		\draw[very thick] (0,0) rectangle (6,6);
		\filldraw[draw=black,fill=white] (2,1) circle (12pt);
		\filldraw[draw=black,fill=white] (2,2) circle (12pt);
		\filldraw[draw=black,fill=white] (0,2) circle (12pt);
		\filldraw[draw=black,fill=white] (1,2) circle (12pt);
		\filldraw[draw=black,fill=white] (3,2) circle (12pt);
		\filldraw[draw=black,fill=white] (4,2) circle (12pt);
		\filldraw[draw=black,fill=white] (4,1) circle (12pt);
		\filldraw[draw=black,fill=white] (3,1) circle (12pt);
		\filldraw[black] (0,1) circle (12pt);
		\filldraw[black] (1,1) circle (12pt);
		\filldraw[black] (1,0) circle (12pt);
		\filldraw[black] (3,1) circle (12pt);
		\filldraw[black] (2,1) circle (12pt);
		\filldraw[black] (5,0) circle (12pt);
		\filldraw[black] (5,1) circle (12pt);
		\fill[draw=black,fill=Red, very thin] (0.25,0.25) rectangle (-0.25,-0.25);
		\fill[draw=black,fill=Blue, very thin] (1.75,0.25) rectangle (2.25,-0.25);
		\fill[draw=black,fill=OliveGreen, very thin] (2.75,0.25) rectangle (3.25,-0.25);
		\fill[draw=black,fill=Blue, very thin] (3.75,0.25) rectangle (4.25,-0.25);
		% Draw the central black dot
	\end{tikzpicture}
	\caption{Tablero con un \textcolor{Red}{ojo} y una \textcolor{Blue}{Unión}.}
	\label{fig:ej_vidas}
\end{wrapfigure}Es crucial diferenciar entre Vivo e Incondicionalmente Vivo. En la Figura~\ref{fig:ej_vidas}, si el negro jugara en el \textcolor{OliveGreen}{punto verde} obtendría vida incondicional; en su estado actual carece de ella. En una partida, un grupo sin vida incondicional puede estar prácticamente vivo, requiriendo apenas una o dos jugadas para asegurarla. La diferencia es justamente esa, la vida incondicional implica que estáticamente algo no pueda morir, la vida en el Go es sobre en cuanto tiempo puedo lograrlo y si los otros jugadores pueden impedirlo.

\subsection{Seki}

\begin{wrapfigure}[10]{r}{0.3\textwidth}
	\vspace{-10pt}
	\centering
	\begin{tikzpicture}[scale=0.485] \tikzstyle{every node}=[font=\small]
		% Draw the background
		\fill[goban] (-0.5,-0.5) rectangle (6.5,6.5);
		
		% Draw the grid lines
		\draw[step=1cm, black] (0,0) grid (6,6);
		
		% Draw the thick border
		\draw[very thick] (0,0) rectangle (6,6);
		\filldraw[draw=black,fill=white] (1,0) circle (12pt);
		\filldraw[draw=black,fill=white] (1,1) circle (12pt);
		\filldraw[draw=black,fill=white] (1,2) circle (12pt);
		\filldraw[draw=black,fill=white] (1,3) circle (12pt);
		\filldraw[draw=black,fill=white] (2,3) circle (12pt);
		\filldraw[draw=black,fill=white] (3,3) circle (12pt);
		\filldraw[black] (4,3) circle (12pt);
		\filldraw[black] (5,3) circle (12pt);
		\filldraw[black] (5,2) circle (12pt);
		\filldraw[black] (5,1) circle (12pt);
		\filldraw[black] (5,0) circle (12pt);
		\filldraw[draw=black,fill=white] (4,0) circle (12pt);
		\filldraw[draw=black,fill=white] (4,1) circle (12pt);
		\filldraw[draw=black,fill=white] (4,2) circle (12pt);
		\filldraw[black] (3,2) circle (12pt);
		\filldraw[black] (2,2) circle (12pt);
		\filldraw[black] (2,1) circle (12pt);
		\filldraw[black] (2,0) circle (12pt);
		
		
		\fill[draw=black,fill=OliveGreen, very thin] (2.75,0.25) rectangle (3.25,-0.25);
		\fill[draw=black,fill=OliveGreen, very thin] (2.75,1.25) rectangle (3.25,0.75);
		% Draw the central black dot
	\end{tikzpicture}
	\caption{Tablero con grupos en seki.}
	\label{fig:seki}
\end{wrapfigure}Habiendo visto la vida incondicional, es tiempo de la vida condicional. Existe otra estructura común en el Go que se denomina seki, esto sucede cuando dos grupos de colores distintos comparten libertades de manera que, si alguno de ellos intenta capturar al otro, el primero termina perdiendo véase Figura~\ref{fig:seki}. En la Figura~\ref{fig:seki}, quien juega primero y toma una de las dos \textcolor{OliveGreen}{libertades} de los grupos del interior, pierde, pues ya no será su turno y el otro jugador puede capturar (y lo hará).
\begin{definition}[Seki sin ojos]
	Sean $I$,$J$ dos grupos de colores distintos. Si existe $R$ una región tal que $|R|\geq 2$ y $L(I)=L(J)=R$.
	Entonces $I, J$ están en seki.
\end{definition}
Esta definición encapsula perfectamente el seki sin ojos pues si algún jugador tratara de jugar en $R$, en algún momento quedará una sola libertad, quien juegue en el turno siguiente concreta la captura. Es necesario pedir que $|R|\geq 2$ de no serlo significa que ambos grupos tienen una sola libertad, basta que quien juega, juegue en esa libertad para capturar. Por el teorema de Benson, un grupo con dos ojos posee vida incondicional. Resta entonces analizar el caso de grupos con un único ojo que se encuentran en seki.

\begin{wrapfigure}[7]{r}{0.4\textwidth}
	\vspace{-10pt}
	\centering
	\begin{tikzpicture}[scale=0.485] \tikzstyle{every node}=[font=\small]
		% Draw the background
		\fill[goban] (-0.5,-0.5) rectangle (6.5,6.5);
		
		% Draw the grid lines
		\draw[step=1cm, black] (0,0) grid (6,6);
		
		% Draw the thick border
		\draw[very thick] (0,0) rectangle (6,6);
		
		\filldraw[black] (1,6) circle (12pt);
		\filldraw[black] (1,5) circle (12pt);
		\filldraw[black] (0,5) circle (12pt);
		\filldraw[black] (2,5) circle (12pt);
		\filldraw[draw=black,fill=white] (3,5) circle (12pt);
		\filldraw[draw=black,fill=white] (3,6) circle (12pt);
		\filldraw[draw=black,fill=white] (4,5) circle (12pt);
		\filldraw[draw=black,fill=white] (5,5) circle (12pt);
		\filldraw[draw=black,fill=white] (5,6) circle (12pt);
		\filldraw[black] (6,6) circle (12pt);
		\filldraw[black] (6,5) circle (12pt);
		\filldraw[black] (6,4) circle (12pt);
		\filldraw[black] (5,4) circle (12pt);
		\filldraw[black] (4,4) circle (12pt);
		\filldraw[black] (3,4) circle (12pt);
		\filldraw[draw=black,fill=white] (2,4) circle (12pt);
		\filldraw[draw=black,fill=white] (1,4) circle (12pt);
		\filldraw[draw=black,fill=white] (0,4) circle (12pt);
		\filldraw[draw=black,fill=white] (2,3) circle (12pt);
		
		
		
		\fill[draw=black,fill=OliveGreen, very thin] (1.75,5.75) rectangle (2.25,6.25);
		\fill[draw=black,fill=Red, very thin] (-0.25,5.75) rectangle (0.25,6.25);
		\fill[draw=black,fill=Red, very thin] (3.75,5.75) rectangle (4.25,6.25);
		
		% Draw the central black dot
	\end{tikzpicture}
	\caption{Tablero con grupos en seki.}
	\label{fig:sekiojos}
\end{wrapfigure}	\begin{definition}[Seki con Ojos]
	Sean $I, J$ dos grupos de piedras de distinto color. 
	Sean $O_I, O_J$ un ojo para $I$ y para $J$ respectivamente.
	Si existe una región $R \subseteq L(I)\cap L(J)$ tal que $L(I)=R \cup O_I$ y $L(J)=R \cup O_J$, entonces $I$ y $J$ están en seki.
\end{definition} Por el Lema~\ref{lem_1_benson}, todo ojo puede reducirse a una única libertad. En consecuencia, ocupar la región compartida R implica depender exclusivamente de dicho ojo, permitiendo al adversario capturar el grupo. Si bien existen variantes más complejas de seki (ojos falsos, doble ko, etc.), estas configuraciones básicas bastan para ilustrar la vida condicional.

\vspace{1cm}
\subsection{Observación sobre el Teorema 2 de Bai}
Presentaremos un contraejemplo al Teorema 2 de Bai \cite{KJB1, KJB2}. Vale destacar que Bai primero se pregunta sobre la inmersión de un grupo incondicionalmente vivo de un tablero bidimensional en un tablero tridimensional. Sin embargo, el Teorema 2 afirma un resultado más fuerte,

«In a 2-person go-game on a 3-dimensional board, every two dimensional x-graph is not safe.»

Dado que Bai fundamenta su análisis en el trabajo de Benson, dentro de nuestro marco teórico este resultado se traduce como:

«En un juego de $Go_{2}^{3}$, ningún grupo de color $x$ de dos dimensiones posee vida incondicional.» 

Para Bai, un «two dimensional x-graph» es un grupo de color $x$ que mantiene fija una de las dimensiones del espacio a una constante. Cabe destacar que, en su formulación, exige un tamaño mínimo de tablero de $9 \times 9 \times 9$. Exhibiremos una configuración que contradice el enunciado del Teorema 2 de Bai.


\begin{wrapfigure}[7]{r}{0.25\textwidth}
	\vspace{-15pt}
	\begin{align*}
		V,N,V,N,V,N,V,N,V
	\end{align*}
	\caption{Capas del contraejemplo.}
	\label{contraej_bai}
\end{wrapfigure}Por la naturaleza tridimensional del tablero, podemos particionarlo a lo largo de su tercera coordenada en distintos «pisos» como tableros de $9\times 9 \times 1$. Consideremos una capa enteramente ocupada por un grupo negro, a la que llamaremos $N$, y una capa enteramente vacía, que llamaremos $V$. Construyendo el tablero de abajo hacia arriba a lo largo de la tercer coordenada, podemos alternar estas capas de la siguiente manera, Figura~\ref{contraej_bai}.

Sea $S$ el conjunto formado por todos los grupos negros de la construcción. Toda intersección de una capa vacía es adyacente a una piedra de alguna de las capas negras vecinas, por lo que la región es saludable. Además, toda frontera de una capa vacía está constituida exclusivamente por piedras pertenecientes a $S$.
Por consiguiente, constituyen al menos dos regiones vitales distintas para cada grupo de $S$. Por el Teorema de Benson (Teorema 1 de Bai), el conjunto $S$ posee vida incondicional. En consecuencia, el Teorema de Benson garantiza la existencia de grupos seguros contenidos en subtableros bidimensionales, contradiciendo el Teorema 2 de Bai. La construcción se generaliza inmediatamente a tableros de cualquier altura. Si el número de capas es impar, basta hacer como se ha mostrado. Si es par, es suficiente reemplazar la primer capa negra por una doble, de dos capas consecutivas de piedras negras, siempre espaciando con capas vacías.

En consecuencia, bajo las definiciones de región vital y vida incondicional adoptadas por Benson y utilizadas por Bai, existen grupos bidimensionales seguros en un tablero tridimensional. Esto contradice el enunciado del Teorema 2.


%\begin{wrapfigure}{r}{0.7\textwidth}
%	Más aún, si omitimos la restricción arbitraria del tamaño del tablero impuesta por Bai, el fenómeno puede ilustrarse en su mínima expresión con la ayuda de la herramienta \href{https://github.com/10500x/3dGo}{3dGo}. 
%	\begin{figure}[H]
	%		\centering
	%		\includegraphics[width=0.8\linewidth]{baicontraej}
	%		\caption{Contraejemplo mínimo al resultado de Bai. Grupo negro de dos dimensiones que posee vida incondicional en un tablero de $2 \times 2 \times 3$.}
	%		\label{fig:baicontraejemplo}
	%	\end{figure}


%
%\vfill

%\section*{Glosario (Parte II)}
%\begin{center}
%	\begin{tabular}{p{2cm} p{4.7cm} p{2cm} p{4.7cm}}
	%		\toprule
	%		\textbf{Símbolo} & \textbf{Significado} & \textbf{Símbolo} & \textbf{Significado} \\
	%		\midrule
	%		$\partial R$ & Frontera de la región $R$ & $V$ & Región vital para un grupo \\
	%		$O$ & Ojo (región vital con un solo grupo frontera) & $U$ & Unión (región vital con más de un grupo frontera) \\
	%		$S,P$ & Conjunto de grupos del mismo color & $\Gamma$ & Juego de Go ($G,E,A,\Delta$) \\
	%		$\Delta$ & Función de transición del estado & $u_t$ & Estado del tablero en el turno $t$ \\
	%		\bottomrule
	%	\end{tabular}
%\end{center}
\clearpage